\documentclass[11pt,a4paper]{ltjsarticle}
\usepackage[margin=2.5cm]{geometry}
\usepackage{amsmath,amssymb,amsthm}
\usepackage{enumitem}
\usepackage[hidelinks]{hyperref}
\usepackage[switch]{lineno}
\setlength{\emergencystretch}{2em}
\newtheorem{theorem}{定理}[section]
\newtheorem{lemma}[theorem]{補題}
\newtheorem{proposition}[theorem]{命題}
\newtheorem{corollary}[theorem]{系}
\newtheorem{definition}[theorem]{定義}
\newcommand{\D}{\mathcal{D}}
\newcommand{\WitnessSet}{\mathsf{W}}
\newcommand{\cf}{\mathrm{cf}}
\newcommand{\lin}{\mathrm{lin}}
\newcommand{\Ccf}[1]{\mathcal C_{#1}^{\cf}}
\newcommand{\Clin}[1]{\mathcal C_{#1}^{\lin}}
\newcommand{\pref}{\mathrm{pref}}
\newcommand{\suf}{\mathrm{suf}}
\newcommand{\keywords}[1]{\smallskip\noindent{\bfseries キーワード:} #1}
\title{固定有限モノイド型付けの下での文脈自由言語の分布学習}
\author{Takayuki Kuriyama\\
独立研究者\\
\texttt{growup.kuriyama@gmail.com}}
\date{}
\begin{document}
\linenumbers
\maketitle

\begin{abstract}
本論文では，有限モノイドへの準同型 $h:\Sigma^*\to M$ の核
$\sim_h$ の下で文脈自由言語の分布学習を研究する．同じ $h$-値をもち，
一つの文脈を共有する非空因子は，すべての文脈を共有することを要求する．
固定 $h$ に対し，集合駆動型再構成作用素 $\mathcal B_h$ は文脈自由
$\sim_h$-代入可能標的を有限個の正例証拠から厳密に再構成し，別の保守的
逐次学習器はすべての非空標的を Gold の意味で極限同定する．再構成と各更新は
多項式時間である．

標準証拠は文法サイズと生成語型付き精密化の厚みに関して多項式であるが，
この型付き厚みは通常の厚みより指数的に大きくなり得る．さらに，一つの固定
二元型付け $h_\circ$ に対して，多項式サイズかつ通常の厚み $O(n)$ の縮約済み
CFG をもつ入れ子状有限標的を構成し，同じ集合駆動型作用素が両標的に特徴標本を
もつなら，大きい方の任意の特徴標本が指数長の語を含むことを示す．一方，
任意の固定接頭辞--接尾辞窓 $(k,\ell)$ では，対応する型付けが Yoshinaka の
$(k,\ell)$-代入可能性をちょうど表し，縮約済み CFG サイズと通常の厚みに関して
多項式の特徴データを得る．任意の固定 $h$ の線形標的では，線形 CFG サイズに
関して多項式の特徴データを得る．さらに，固定窓階層からの正則・非正則線形・
非線形の分離例を与え，一種類の括弧の Dyck 言語が可認識代入可能族の外にある
ことを示す．
\end{abstract}

\section{序論}

\subsection{背景と動機}

Gold~\cite{gold1967} は，すべての有限言語と少なくとも一つの無限言語を含む言語クラスは，正例データのみから極限同定できないことを示した．この否定的結果は，文法推論における広範な研究の流れを生み出した．一般的背景については，たとえば de~la~Higuera~\cite{higuera2010} を参照されたい．文脈自由言語について正例学習可能性を回復する標準的な方法の一つが，\emph{分布的}アプローチである．そこでは，異なる文脈間で比較してよい部分文字列を制限することで学習可能性を得る．古典的な例として，Clark--Eyraud の代入可能性~\cite{clark-eyraud2007} および Yoshinaka の $(k,\ell)$-代入可能性~\cite{yoshinaka2008} がある．より一般には，適切な有限制御の仮定の下で，他の文法クラスに対しても正例データからの学習可能性が示されている．たとえば $k$-valued 範疇文法の正例学習可能性については Kanazawa~\cite{kanazawa1996} を参照されたい．

本論文では，この分布的研究計画の可認識合同関係版を研究する．$h:\Sigma^*\to M$ を有限モノイドへの準同型とし，$\sim_h:=\ker h$，すなわち $x\sim_h y \Longleftrightarrow h(x)=h(y)$ とおく．$x\in\Sigma^*$ に対し，
\[
\D_L(x):=\{(u,v)\in\Sigma^*\times\Sigma^*:uxv\in L\}
\]
を $L$ における $x$ の分布とする．言語 $L\subseteq\Sigma^*$ が \emph{$\sim_h$-代入可能}であるとは，すべての $x,y\in\Sigma^+$
について%
\footnote{Clark--Eyraud では $a^+$ が代入可能として扱われるが，$\lambda$ 自身を比較対象に含めると $a^+$ は代入可能でなくなる．この点については Clark~\cite{clark2013} も参照されたい．}，$h(x)=h(y)$ かつ $\D_L(x)\cap\D_L(y)\neq\emptyset$ ならば $\D_L(x)=\D_L(y)$ が成り立つことをいう．したがって分布の完全な一致を要求するのは，あらかじめ固定された有限代数的要約の下ですでに一致する文字列同士に限られる．

ここで，可認識合同関係は自然な比較機構を与える．ある有限モノイド準同型 $h:\Sigma^*\to M$ と $\mathsf{Acc}\subseteq M$ に対して $L=h^{-1}(\mathsf{Acc})$ であるならば，同じ $h$-値をもつ任意の二つの文字列は自動的に同じ継続挙動をもち，したがって同じ分布をもつ．特に，すべての正則言語は，\emph{ある}有限モノイド準同型に対して $\sim_h$-代入可能な言語からなる，より広いクラス $\mathsf{RS}$（第~\ref{sec:framework}節で形式的に定義する）に属する．この観察は，有限モノイド型付けが有界文脈の場当たり的な一般化ではなく，有限な比較バイアスを与える標準的な代数的源であることを示している．

本論文の焦点は \emph{固定-$h$} の状況である．同じ有限モノイド準同型 $h:\Sigma^*\to M$ が，文脈自由 $\sim_h$-代入可能言語のクラス $\Ccf{h}$ 全体にわたって一様に固定され，各標的についてデータ観測後に別々に選ばれるのではない．したがって固定 $h$ は，Yoshinaka の階層における固定された組 $(k,\ell)$ と同じクラスレベルの役割を果たす．有用な型付けを選ぶ問題と，型付けがすでに指定された後に学習する問題とを分離する点で，この観点はオートマトン推論における，外部から与えられる型付け・領域背景知識の利用~\cite{coste-typing2004} と方法論的に近い．Coste らは専門家の型付けを有限状態の型オートマトンで表現し，ある記号出現の型がその左右文脈に依存しうることを許したうえで，推論されたオートマトンに埋め込まれる型付けが専門家の型付けと両立するよう状態併合を制約する．これに対して本論文の $h$ は，オートマトン構造中の出現位置に文脈依存の型を与えるのではなく，終端部分文字列そのものに合成的な有限型を与える点で形式的には異なるが，固定された有限比較バイアスとしては類似した役割を果たす．

制御集合アプローチとの対比もある．Takada の 1988 年の結果は，固定された普遍 even-linear 文法上の正則制御集合の学習へ even-linear 言語の推論を還元する~\cite{takada1988}．1995 年の総合的な章ではこの構図がさらに明示され，固定された普遍かつ非曖昧な even-linear 文法上で各標的は一意な標準正則制御集合をもち，even-linear 言語の学習が正則言語の学習へ還元される．同じ章では，線形文法および equal-matrix 文法についても正則制御集合による表現がより広く展開されている~\cite{takada1995}．これらの枠組みでは，標的ごとの正則制御言語が生成規則ラベル列を制約する．これに対し，本論文の準同型は終端生成語に作用し，クラス全体で固定された有限な \emph{比較バイアス} を与える．すなわち，$h$ を固定しても一つの標的言語が決まるのではなく，どの部分文字列を比較してよいかが制限されるだけである．制御集合に関する正例学習では，正則制御集合にさらに制限を課すことで学習可能な部分クラスが得られる~\cite{koshiba-makinen-takada1997}．ここでは位置づけのためにこの対比を用いるだけであり，control-set 方法論全体との同値性・包含関係・クラス分離はいずれも主張しない．

\subsection{貢献と主要結果}

本論文の主要な貢献は三つである．

\begin{enumerate}[label=\textup{(\arabic*)},leftmargin=2.4em]
\item \emph{固定 $h$ に対する学習．}
集合駆動型作用素 $\mathcal B_h$ は，$\Ccf{h}$ の任意の標的を有限個の
標準正例証拠から厳密に再構成する．別の保守的逐次学習器 $\mathcal A_h$ は
すべての非空標的を Gold の意味で同定し，再構成と各更新は蓄積標本に関して
多項式時間である．

\item \emph{定量的境界．}
任意の固定 $h$ では，標準証拠は文法サイズと生成語型付き精密化の厚みに
関して多項式であるが，通常の厚みが $1$ でも型付き厚みは指数的になり得る．
一つの固定二元型付け $h_\circ$ に対しては，
$\Ccf{h_\circ}$ 全体に特徴標本を与える集合駆動型作用素には，縮約済み CFG
サイズと通常の厚みによる一般の多項式データ上界が存在しない．これに対し，
任意の固定窓 $(k,\ell)$ では Yoshinaka の $(k,\ell)$-代入可能性とちょうど
対応する多項式厚み上界が得られ，任意の固定 $h$ の線形標的では線形文法サイズ
に関して多項式の特徴データが得られる．

\item \emph{表現力上の位置．}
適切な有限アルファベット上で，可認識代入可能族は固定窓階層を正則，
非正則線形，非線形文脈自由の各水準で真に拡張する．一方，一種類の括弧の
Dyck 言語は $\mathsf{RS}$ の外にあり，
$\mathsf{RS}\cap\mathrm{CFL}$ は Clark の congruential family に真に含まれる．
\end{enumerate}

\subsection{論文の構成}
第~\ref{sec:prelim}--\ref{sec:framework}節で記法と古典的特殊場合を整え，
第~\ref{sec:algo}--\ref{sec:complexity}節で学習結果を証明する．
第~\ref{sec:typed-thickness}--\ref{sec:linear}節で定量的境界と二つの正の場合を扱い，
第~\ref{sec:boundaries}節で表現力上の境界を与える．長い正規化証明は付録に置く．

\section{準備}
\label{sec:prelim}

$\Sigma$ を有限アルファベットとし，$\Sigma^*$ を単位元 $\lambda$ をもつ自由モノイドとする．また $\Sigma^+:=\Sigma^*\setminus\{\lambda\}$ とおく．対 $(u,v)\in\Sigma^*\times\Sigma^*$ を \emph{終端文脈} と呼ぶ．$L\subseteq\Sigma^*$ と $x\in\Sigma^*$ に対し，$L$ における $x$ の \emph{分布} を
\[
\D_L(x):=\{(u,v)\in\Sigma^*\times\Sigma^*:uxv\in L\}
\]
と定める．$L$ の \emph{構文合同関係} は，$x,y\in\Sigma^*$ に対して $x\equiv_L y \Longleftrightarrow \D_L(x)=\D_L(y)$ により定める．$x\equiv_L y$ のとき，$x$ と $y$ は \emph{構文合同} であるといい，$x$ の $\equiv_L$-同値類を $[x]_{\equiv_L}$ と書く．

文脈自由文法（CFG）は有限非終端記号集合 $V$，開始記号 $S\in V$，および $A\in V$, $\alpha\in(V\cup\Sigma)^*$ に対する生成規則 $A\to\alpha$ からなる組 $G=(V,\Sigma,P,S)$ である．\emph{文形式}とは $(V\cup\Sigma)^*$ の任意の語である．$G$ における一段階導出を $\alpha\Rightarrow_G\beta$ と書き，その反射推移閉包を $\Rightarrow_G^*$ と書く．したがって
\[
L(G):=\{w\in\Sigma^*:S\Rightarrow_G^*w\}.
\]
$A\in V$ に対して
\[
L_G(A):=\{w\in\Sigma^*:A\Rightarrow_G^*w\}
\]
とおく．非終端記号 $X\in V$ に対し，$S\Rightarrow_G^*uXv$ が成り立つとき，終端文脈 $(u,v)$ を $G$ における $X$ の \emph{終端到達文脈} と呼ぶ．非終端記号が $S$ から到達可能な文形式中に現れるとき \emph{到達可能} であるといい，終端語を導出するとき \emph{生産的} であるという．到達可能性または生産性の少なくとも一方を欠く非終端記号を \emph{無用} といい，$\lambda$ を導出する非終端記号を \emph{空語導出可能（nullable）} という．$A,B\in V$ に対する生成規則 $A\to B$ を \emph{単位規則} という．\emph{単位閉包}とは，単位生成規則が非終端記号上に誘導する関係の反射推移閉包である．すべての非終端記号が到達可能かつ生産的である文法を \emph{縮約済み} という．非空言語を生成する文法について，その \emph{トリム}とは，すべての無用な非終端記号と，それらに接続する生成規則を削除して得られる部分文法をいう．各右辺に高々一つの非終端記号しか現れない CFG を \emph{線形} という．\emph{二分化}とは，新たな非終端記号を用いて長い右辺を同値な二項生成規則の鎖へ置き換えることをいう．
$\mathrm{CFL}$ を文脈自由言語全体のクラスとする．任意の有限表現 $r$ に対し，$|r|$ は固定した妥当な符号化の長さを表すものとする．文法については，これは通常の生成規則中の記号数と多項式の範囲で同値である．Yoshinaka の固定窓解析~\cite{yoshinaka2008} に従い，縮約済み文法 $G=(V,\Sigma,P,S)$ について\emph{厚み}パラメータを
\[
 \tau_G:=\max_{A\in V}\min\{\,|w|:A\Rightarrow_G^*w,\ w\in\Sigma^*\,\}
\]
と定める．すなわち $\tau_G$ は，各非終端記号の最短終端生成語の長さを取り，その最大値をとったものである．

\subsection{正例データからの学習}\label{sec:learning}

Gold の explanatory model~\cite{gold1967} を用いる．$H_0$ を，一つの開始記号をもち生成規則をもたない固定 CFG とし，$L(H_0)=\emptyset$ とする．以下のテキスト規約は非空の標的言語に限定する．空言語はこのテキストに基づく Gold の記述から除外し，別途 $H_0$ により表す．同値に，通常の pause 記号をテキストに加えることで空標的を含めることもできるが，余分な記法を避けるためにここでは非空テキスト規約を用いる．非空言語 $L$ に対し，\emph{テキスト}とは，$L$ の各語が少なくとも一度は現れる $L$ 上の無限列 $T=(w_1,w_2,\ldots)$ である．$T[n]=(w_1,\ldots,w_n)$ および $K_n:=\{w_1,\ldots,w_n\}$ と書く．学習器 $\mathcal A$ が非空言語 $L$ を \emph{極限同定}するとは，$L$ の任意のテキスト $T$ に対して，ある $n_0$ と文法 $H$ が存在し，すべての $n\ge n_0$ について $\mathcal A(T[n])=H$ かつ $L(H)=L$ が成り立つことをいう．したがって，仮説表現そのものが最終的に安定する．逐次学習器が \emph{保守的} であるとは，新たに提示された正例データが現在の仮説によって生成されない場合にのみ，その現在の仮説を変更することをいう．

この逐次学習器とは別に，集合駆動型の有限標本作用素 $\mathcal B$ を考える．ここで $\mathcal B(K)$ は有限集合 $K$ のみに依存し，提示順序や重複回数には依存しない．任意の有限標本 $K$ に対して $K\subseteq L(\mathcal B(K))$ が成り立つとき，$\mathcal B$ は \emph{標本整合的} であるという．有限集合 $C\subseteq L$ が $\mathcal B$ に対する \emph{特徴標本} であるとは，任意の有限 $K$ について $C\subseteq K\subseteq L$ ならば $L(\mathcal B(K))=L$ が成り立つことをいう．これは de~la~Higuera の多項式文法推論枠組み~\cite{higuera1997} における有限標本部分である．標本サイズを
\[
\|K\|:=\sum_{w\in K}(|w|+1)
\]
で測る\footnote{$+1$ は，$K=\{\lambda\}$ のような非空標本のサイズが $0$ になることを防ぐ．また $|K|\le \|K\|$ となるため，正例数も同じサイズ尺度で抑えられる．}．したがって $\lambda$ も 1 を寄与する．以下の定量的結果では，$\mathcal B(K)$ の多項式時間構成と，多項式サイズの特徴データを組み合わせる．

\subsection{固定有限モノイド型付け}\label{sec:fixed-typing}

有限モノイド $M$ と準同型 $h:\Sigma^*\to M$ を固定する．これらは $M$ の乗法表および各 $a\in\Sigma$ に対する文字像 $h(a)$ によって表現される．したがって $h(w)$ は左から右への積によって $O(|w|)$ 時間で計算できる．以下で固定準同型が \emph{効率的に計算可能} であるというのはこの意味である．

\begin{definition}[固定-$h$ 代入可能性]\label{def:hsubst}
言語 $L\subseteq\Sigma^*$ が \emph{$\sim_h$-代入可能} であるとは，すべての $x,y\in\Sigma^+$ について，$h(x)=h(y)$ と $\D_L(x)\cap\D_L(y)\ne\varnothing$ から $\D_L(x)=\D_L(y)$ が従うことをいう．
\end{definition}
空語は内部の代入可能因子としては用いない．$\lambda\in L$ の場合には，再構成において開始記号を通じて別途処理する．

非空内部因子 $x$ に対して，$h(x)$ をその \emph{生成語型} と呼ぶ．文脈自由 $\sim_h$-代入可能言語のクラスを $\Ccf{h}$ と書き，
\[
\Clin{h}:=\{L\in\Ccf{h}:L\text{ はある線形 CFG により生成される}\}
\]
と定める．したがって $\Clin{h}$ は線形言語からなる部分クラスである．以下の多項式データ上界は，標的について選ばれた線形 CFG 表現を基準に測る．後で構成される再構成仮説自体は線形文法である必要はない．

\section{固定可認識代入可能性}
\label{sec:framework}

固定準同型 $h$ に対し，
\[
\mathsf{RS}_h:=\{L\subseteq\Sigma^*:L\text{ は $\sim_h$-代入可能}\}
\]
とおく．したがって $\Ccf{h}=\mathsf{RS}_h\cap\mathrm{CFL}$ であり，$\Clin{h}$ は $\mathsf{RS}_h\cap\mathrm{CFL}$ に属し，かつある線形 CFG により生成される言語からなる．表現力の比較のためだけに，$h$ が定義域 $\Sigma^*$ をもつ有限モノイド準同型全体を動くものとして，
\[
\mathsf{RS}:=\bigcup_h\mathsf{RS}_h
\]
とする．このようなクラス記法はすべて周囲のアルファベットに相対的である．アルファベットを明示する必要があるときは，$\Gamma^*$ 上の対応するクラスを $\mathsf{RS}_h(\Gamma)$ および $\mathsf{RS}(\Gamma)$ と書く．クラス包含は常にアルファベットごとに理解する．本論文の学習結果は一つの固定 $h$ を対象とする．非空アルファベット上では，$\mathsf{RS}$ も $\mathsf{RS}\cap\mathrm{CFL}$ も正例データから極限同定できない．\footnote{命題~\ref{prop:regular-auto} で示すように，すべての正則言語は $\mathsf{RS}$ に属する．したがって，非空アルファベット $\Sigma$ 上では，$\mathsf{RS}$ も $\mathsf{RS}\cap\mathrm{CFL}$ もすべての有限言語を含み，さらに無限正則言語 $\Sigma^*$ も含む．ゆえに両クラスは superfinite であり，Gold の定理~\cite{gold1967} により正例データから極限同定できない．}ただし，型付けバイアスの精密化について単純な単調性がある．$h':\Sigma^*\to M'$ が $\ker h'\subseteq\ker h$ という意味で $h$ を精密化するとき，$\mathsf{RS}_h\subseteq\mathsf{RS}_{h'}$ が成り立つ．実際，$h'$ の下で等しいことは $h$ の下で等しいことを意味するので，$h'$ が要求するすべての比較はすでに $h$-代入可能性によって覆われている．したがって，より細かい型付けは比較すべき因子対を減らすことにより，より大きな固定-$h$ クラスを許す．ただし，この利得は以下の量的上界において無償ではない．型付けを精密化すると，後述する型付き非終端記号数が増え，その結果，証拠集合サイズも増える可能性がある．特に，二つの有限型付け $h,g$ に対して積準同型 $h\times g$ は $\ker(h\times g)=\ker h\cap\ker g$ をもち，両者を精密化する．したがって $\mathsf{RS}_h\cup\mathsf{RS}_g\subseteq\mathsf{RS}_{h\times g}$ であり，窓情報と独立な有限の領域固有特徴を一つの固定型付けに合成できる．

型付けに関するもう一つの有用な観察として，単純な十分条件が得られる．$\ker h$ を $\Sigma^+$ に制限した関係が $\equiv_L$ を精密化するならば\footnote{もし$\ker h\subseteq\equiv_L$ を $\Sigma^*$ 全体で要求してしまうと，例えば自明モノイドでは $\lambda$ とすべての非空語が同じ型になるため，通常は代入可能とされる $a^+$ も $\lambda\not\equiv_{a^+} a$ により除外されてしまう．}，$L\in\mathsf{RS}_h$ である．なぜなら，$h$-型が等しい語はすでに構文合同だからである．しかしこれは必要条件ではない．後の命題~\ref{prop:linear-separator-example} は，四元の固定-$h$ に属しながら構文合同関係が無限指数をもつ例を与える．任意の有限モノイド核の指数は有限なので，固定-$h$ 代入可能性は，$\ker h$ が完全な構文合同関係を精密化することよりも真に弱い．

\begin{proposition}\label{prop:regular-auto}
言語が正則であることと，ある有限モノイド $M$ への準同型 $h:\Sigma^*\to M$ およびある $\mathsf{Acc}\subseteq M$ に対して $h^{-1}(\mathsf{Acc})$ の形に書けることは同値である．さらに，任意の固定 $h$ と $\mathsf{Acc}\subseteq M$ に対し，言語 $h^{-1}(\mathsf{Acc})$ は $\sim_h$-代入可能である．したがって，すべての正則言語は $\mathsf{RS}$ に属する．
\end{proposition}

\begin{proof}
最初の主張は，有限モノイドによる正則言語の標準的特徴づけである~\cite{pin2025}．
$L=h^{-1}(\mathsf{Acc})$ かつ $h(x)=h(y)$ ならば，任意の $u,v$ に対して
$h(uxv)=h(uyv)$ である．したがって $\D_L(x)=\D_L(y)$ である．
\end{proof}

\subsection{古典的特殊場合}

言語 $L\subseteq\Sigma^*$ が \emph{Clark--Eyraud 代入可能} であるとは，すべての $x,y\in\Sigma^+$ について，$\D_L(x)\cap\D_L(y)\ne\varnothing$ ならば $\D_L(x)=\D_L(y)$ が成り立つことをいう．定義~\ref{def:hsubst} の非空因子規約の下では，これはちょうど自明モノイドの場合である．特に，Yoshinaka の $(0,0)$ の場合に一致する．同じ運用上の規約については Clark~\cite{clark2013} も参照されたい．

\begin{proposition}[Clark--Eyraud の特殊場合]\label{prop:ce-special}
すべての Clark--Eyraud 代入可能言語は $\mathsf{RS}$ に属する．
\end{proposition}
\begin{proof}
一元モノイドをとる．すると型条件は自明となり，定義~\ref{def:hsubst} は通常の代入可能性そのものになる．
\end{proof}

$k,\ell\ge0$ に対し，$\pref_k(w)$ および $\suf_\ell(w)$ を，それぞれ存在するときの指定長の接頭辞と接尾辞とし，$\pref_0(w)=\suf_0(w)=\lambda$ とする．

Yoshinaka の元の定義~\cite{yoshinaka2008} では，長さ $k$ の左窓を $v\in\Sigma^k$，長さ $\ell$ の右窓を $u\in\Sigma^\ell$ と書く．以下では一般の文脈記法 $(u,v)$ との衝突を避けるため，この二つの窓変数だけを $p,q$ と改名する．したがって，言語 $L$ が \emph{$(k,\ell)$-代入可能} であるとは，すべての $p\in\Sigma^k$, $q\in\Sigma^\ell$ および $x_1,x_2,y_1,y_2,z_1,z_2\in\Sigma^*$ について，$py_1q,py_2q\ne\lambda$ かつ
\[
\begin{split}
&x_1py_1qz_1\in L,\quad x_1py_2qz_1\in L,\quad
x_2py_1qz_2\in L\\
&\hspace{35mm}\Longrightarrow\quad x_2py_2qz_2\in L
\end{split}
\]
が成り立つことをいう．

\begin{proposition}\label{prop:yl-special}
任意の固定 $k,\ell\ge0$ に対して有限モノイド準同型 $h_{k,\ell}$ が存在し，任意の言語 $L\subseteq\Sigma^*$ について，$L$ が $(k,\ell)$-代入可能であることと，$L$ が $\sim_{h_{k,\ell}}$-代入可能であることは同値である．
\end{proposition}

\begin{proof}

$m_0:=\max\{1,k+\ell\}$ とおく．$|w|<m_0$ のとき $h_{k,\ell}(w)=w$ とし，それ以外では
\[
h_{k,\ell}(w)=\langle\pref_k(w),\suf_\ell(w)\rangle
\]
とする．1 ビットのタグを用いて，短い語を表す要素と長い語を表す要素とを互いに素に区別する．さらに $M_{k,\ell}:=h_{k,\ell}(\Sigma^*)$ とおく．短い語 $x,y$ と像中のタグ付き要素に対し，積を
\[
\begin{aligned}
 x\cdot y&:=h_{k,\ell}(xy), &
 x\cdot\langle p,q\rangle&:=\langle\pref_k(xp),q\rangle,\\
 \langle p,q\rangle\cdot y&:=\langle p,\suf_\ell(qy)\rangle, &
 \langle p,q\rangle\cdot\langle p',q'\rangle&:=\langle p,q'\rangle
\end{aligned}
\]
と定める．四つの場合を直接確認すれば $h_{k,\ell}(xy)=h_{k,\ell}(x)\cdot h_{k,\ell}(y)$ が得られる．$h_{k,\ell}$ は全射なので，結合律は連接の結合律から従い，$h_{k,\ell}(\lambda)$ が単位元である．したがって $M_{k,\ell}$ は有限モノイドであり，$h_{k,\ell}$ は準同型である．

$s,t\in\Sigma^+$ かつ $h_{k,\ell}(s)=h_{k,\ell}(t)$ かつ $\D_L(s)\cap\D_L(t)\ne\varnothing$ と仮定する．$|s|<m_0$ ならば $s=t$ である．そうでなければ $s,t$ はともに長く，同じ長さ $k$ の接頭辞 $p$ と長さ $\ell$ の接尾辞 $q$ をもつので，$s=py_1q$，$t=py_2q$ と書ける．$(u_0,v_0)$ を $s,t$ の共通文脈とし，$(u',v')$ を $s$ の任意の文脈とする．三つの語
\[
u_0py_1qv_0,\quad u_0py_2qv_0,\quad u'py_1qv'
\]
は $L$ に属する．$(k,\ell)$-代入可能性を適用すると $u'py_2qv'\in L$ を得る．よって $\D_L(s)\subseteq\D_L(t)$ であり，対称性から等号が成り立つ．

逆に，$L$ が $\sim_{h_{k,\ell}}$-代入可能であるとし，$(k,\ell)$-代入可能性の前提を満たす文字列をとる．$s:=py_1q$, $t:=py_2q$ とおくと，定義により両者は非空である．$k+\ell>0$ なら $|s|,|t|\ge k+\ell=m_0$ であり，両者の型はともに $\langle p,q\rangle$ である．$k=\ell=0$ なら，非空性から $|s|,|t|\ge1=m_0$ であり，両者は唯一の非空語型 $\langle\lambda,\lambda\rangle$ をもつ．よって，どちらの場合でも $h_{k,\ell}(s)=h_{k,\ell}(t)$ である．最初の二つの前提語から $s,t$ は文脈 $(x_1,z_1)$ を共有し，第三の前提語から $(x_2,z_2)\in\D_L(s)$ である．固定-$h_{k,\ell}$ 代入可能性により $(x_2,z_2)\in\D_L(t)$ が従い，これは必要な第四の語そのものである．
\end{proof}

\phantomsection\label{crit:kl}
\paragraph{固定窓代入可能性の反例基準．}
$k,\ell\ge0$ とし，$m_0:=\max\{1,k+\ell\}$ とおく．
$s,t\in\Sigma^*$ が $|s|,|t|\ge m_0$，
$\pref_k(s)=\pref_k(t)$，$\suf_\ell(s)=\suf_\ell(t)$，
$\D_L(s)\cap\D_L(t)\ne\emptyset$ を満たし，かつ
$\D_L(s)\ne\D_L(t)$ であるならば，$L$ は $(k,\ell)$-代入可能ではない．
実際，長さの仮定と接頭辞--接尾辞の一致から
$h_{k,\ell}(s)=h_{k,\ell}(t)$ である一方，
$s,t$ は共通文脈をもちながら分布が異なるので，
$L$ は $\sim_{h_{k,\ell}}$-代入可能ではない．
命題~\ref{prop:yl-special} より結論が従う．

したがって，固定 $(k,\ell)$ クラスはちょうど $\mathsf{RS}_{h_{k,\ell}}$ である．第~\ref{sec:fixed-window-thick}節では，この特殊場合に本論文の再構成作用素が Yoshinaka の文法サイズと厚みに関する有限標本尺度を回収することを確認する．ここで $k,\ell$ は固定されており，窓パラメータを可変とした一様な多項式上界は主張しない．また，厚み付きデータの主張は集合駆動型作用素 $\mathcal B_{h_{k,\ell}}$ による一括再構成に関するものであり，保守的逐次学習器の最終的安定化は別に Gold の極限同定として示す．$k=\ell=0$ のとき，準同型 $h_{0,0}$ は $\lambda$ と非空語を区別し，非空内部因子だけを比較する定義~\ref{def:hsubst} の下では一元自明モノイドと同じ代入可能性条件を誘導する．アルファベット $\Gamma$ に対し，
\[
\mathsf{KL}(\Gamma):=
\bigcup_{k,\ell\ge0}
\{L\subseteq\Gamma^*:L\text{ は $(k,\ell)$-代入可能}\}
\]
と定め，現在のアルファベット上では $\mathsf{KL}:=\mathsf{KL}(\Sigma)$ と略記する．

\subsection{主定理}

\begin{theorem}[固定-$h$ 学習定理]\label{thm:main}
有限モノイドへの効率的に計算可能な準同型 $h:\Sigma^*\to M$ を固定する．
集合駆動型再構成作用素 $\mathcal B_h$ と保守的逐次学習器 $\mathcal A_h$ が
存在して，次を満たす．
\begin{enumerate}[label=\textup{(\roman*)},nosep]
\item $\mathcal B_h$ は標本整合的で，$\mathcal B_h(K)$ は
$\|K\|$ に関して多項式時間で構成できる．
\item すべての $L\in\Ccf{h}$ は $\mathcal B_h$ に対する有限特徴標本をもつ．
$L\ne\emptyset$ で $G$ が $L$ の縮約済み SSBNF 文法なら，その一つは
$|G|$ と $\tau_h^{\mathrm{typ}}(G)$ の多項式サイズに選べる．
\item $\mathcal A_h$ はすべての非空 $L\in\Ccf{h}$ を Gold の意味で同定し，
各更新はそれまでに見たデータのサイズに関して多項式時間である．
\item 固定 $(k,\ell)$ に対し，$\Ccf{h_{k,\ell}}$ の標的を表す任意の縮約済み
CFG $R$ について，(ii) の標本は $|R|$ と $\tau_R$ の多項式に選べる．
また $L\in\Clin{h}$ なら，任意の線形 CFG $R$ に対して $|R|$ の多項式に選べる．
\end{enumerate}
これに対し，ある固定二元準同型 $h_\circ$ について，
$\Ccf{h_\circ}$ の入れ子状有限標的が存在し，両者は多項式サイズ，
通常の厚み $O(n)$ の縮約済み CFG をもつが，両標的に特徴標本をもつ任意の
集合駆動型作用素に対し，大きい方の標的の任意の特徴標本サイズは
$2^{\Omega(n)}$ である．特に $\mathcal B_{h_\circ}$ にもこの下界が適用される．
\end{theorem}

Gold の superfinite obstruction との矛盾はない．相異なる非空語 $x,y$ を
$h(x)=h(y)$ となるように選び，$|r|>\max\{|x|,|y|\}$ とすれば，
$x,y$ は $\{x,y,rx\}$ で空文脈を共有する一方，$(r,\lambda)$ は $x$ の
文脈にしかならない．従ってこの有限言語は $\sim_h$-代入可能ではない．

\section{再構成作用素と逐次学習器}
\label{sec:algo}

\subsection{有限標本からの再構成}
有限標本 $K\subseteq\Sigma^*$ に対し，
\[
\widehat V(K):=\{[x:u,v]:uxv\in K,\ x\in\Sigma^+\}
\]
とおく．これらは構文的三つ組である．すなわち規約として，$[x:u,v]=[x':u',v'] \Longleftrightarrow x=x',\ u=u',\ v=v'$ とする．特に，同じ標本語から生じる場合であっても，因子文字列または文脈文字列が異なる二つの出現は，異なる仮説非終端記号である．文法 $\mathcal B_h(K)=\widehat G(K)$ は，非終端記号集合 $\widehat V(K)\cup\{\widehat S\}$ と開始記号 $\widehat S$ をもつ．表示されるすべての角括弧付き非終端記号が $\widehat V(K)$ に属する場合に，次の生成規則をもつ．
\[
\begin{array}{rll}
(R1)&[xy:u,v]\to[x:u,yv][y:ux,v],&x,y\in\Sigma^+,\\[1mm]
(R2)&[x:u,v]\to[x:u',v'],&uxv,u'xv'\in K,\\[1mm]
(R3)&[x:u,v]\to[x':u,v],&h(x)=h(x'),\\[1mm]
(R4)&[a:u,v]\to a,&a\in\Sigma,\\[1mm]
(R5)&\widehat S\to[w:\lambda,\lambda],&w\in K\cap\Sigma^+.
\end{array}
\]
$\lambda\in K$ のときは，規則~(R5) に $\widehat S\to\lambda$ も含める．直観的には，$[x:u,v]$ は観測された因子 $x$ を文脈 $(u,v)$ の中で置き換える候補を表す．$\sim_h$-代入可能な標的に対しては，以下の健全性不変量により，導出されるすべての語が $x$ と同じ $h$-型をもち，分布的に交換可能であることが保証される．規則~(R2) は文脈輸送規則である．同じ観測因子は，その因子自身が観測された任意の文脈で表現できる．二つの周囲文脈の $h$-型に関する条件は不要である．規則~(R3) は，一つの固定生成語型の内部で代表元を変更する．その条件 $h(x)=h(x')$ は，健全性不変量の生成語型部分を保存するためにまさに用いられる．その不変量はその後規則~(R1) および~(R2) で代入可能性を適用するときに用いられる．

\begin{lemma}[標本整合性]\label{lem:sample-consistency}
任意の有限正例標本 $K$ に対して，$K\subseteq L(\mathcal B_h(K))$ が成り立つ．
\end{lemma}
\begin{proof}
$w=a_1\cdots a_m\in K\cap\Sigma^+$ とする．規則~(R5) により $[w:\lambda,\lambda]$ に入る．この $w$ の出現を任意に二分木分解し，それに沿って規則~(R1) を繰り返し適用すると，一文字因子へ到達する．中間で必要となる観測非終端記号はすべて，同じ標本語の因子と，それによって誘導される文脈から生じるため存在する．最後に規則~(R4) が $a_1\cdots a_m$ を出力する．$\lambda\in K$ の場合は，開始規則がこれを直接出力する．
\end{proof}

また $K\subseteq K'$ ならば，$K$ から作られたすべての観測非終端記号と規則~(R1)--(R5) のすべてのインスタンスは $K'$ に対しても残る．したがって $\widehat G(K)$ は $\widehat G(K')$ の部分文法であり，
\[
L(\mathcal B_h(K))\subseteq L(\mathcal B_h(K')).
\]
ゆえに保守的学習器が再構成を行うとき，仮説言語は拡大することしかない．

\begin{theorem}[健全性]\label{thm:soundness}
$K\subseteq L$ かつ $L$ が $\sim_h$-代入可能ならば，$L(\widehat G(K))\subseteq L$ である．
\end{theorem}

\begin{proof}
$[x:u,v]$ を根とする導出木の高さについて帰納法を行い，次を同時に証明する．
\begin{equation}\label{eq:soundness-invariant}
 [x:u,v]\Rightarrow_{\widehat G(K)}^*w
 \quad\Longrightarrow\quad
 uwv\in L\ \text{かつ}\ h(w)=h(x).
\end{equation}
規則~(R1)--(R4) は消去生成規則を含まず，すべての語彙規則は一つの終端記号を出力する．したがって，開始記号以外の仮説非終端記号から導出されるすべての終端語は非空である．特に，以下で定義~\ref{def:hsubst} を適用する語 $w,w_1,w_2$ は $\Sigma^+$ に属する．規則~(R4) の場合は明らかである．規則~(R3) については，帰納法の仮定から $h(w)=h(x')=h(x)$ が得られ，またすでに $uwv\in L$ が得られる．規則~(R2) については，帰納法の仮定から $u'wv'\in L$ および $h(w)=h(x)$ が得られる．さらに $[x:u',v']$ と $[x:u,v]$ は $\widehat V(K)$ に属するので，その定義から $u'xv',uxv\in K\subseteq L$ である．したがって $x$ と $w$ は文脈 $(u',v')$ を共有し，代入可能性により $\D_L(x)=\D_L(w)$ を得る．よって $uwv\in L$ である．

規則~(R1) について，導出語を $w_1w_2$ と書く．二つの帰納法の仮定から
\[
 uw_1yv\in L,\qquad uxw_2v\in L,\qquad
 h(w_1)=h(x),\qquad h(w_2)=h(y)
\]
が得られる．$uxyv\in K$ なので，$x$ と $w_1$ は文脈 $(u,yv)$ を共有し，同じ型をもつ．したがって両者の分布は一致する．第二の帰納法の仮定により $(u,w_2v)\in\D_L(x)$ である．この文脈に分布の等しさを適用すると $uw_1w_2v\in L$ を得る．さらに帰納法の仮定から $h(w_1w_2)=h(xy)$ である．これで~\eqref{eq:soundness-invariant} が証明された．開始規則を用いれば $L(\widehat G(K))$ に対する主張が従う．$\lambda$ は $\lambda\in K\subseteq L$ の場合にのみ存在する．
\end{proof}

\subsection{逐次学習器}

$\mathcal B_h$ を各データ到着時に再計算しても，文法表現は一般には安定しない．
\footnote{たとえば $L=\{a^ncb^n:n\ge0\}$ では，長い正例が加わるたびに
$[a^ncb^n:\lambda,\lambda]$ のような新しい三つ組が増える．
set-based と incremental learning の差については
Eyraud--Heinz--Yoshinaka~\cite[Appendix]{eyraud-heinz-yoshinaka2016} も参照されたい．}
そこで Clark--Eyraud の SGL および Yoshinaka の $(k,\ell)$-SGL と同じ
保守的ラッパーを用いる~\cite{clark-eyraud2007,yoshinaka2008}．固定空文法
$H_0$ から始め，$w_n$ を読んだ後
\[
 H_n:=
 \begin{cases}
 H_{n-1},&w_n\in L(H_{n-1}),\\
 \mathcal B_h(K_n),&w_n\notin L(H_{n-1}),
 \end{cases}
 \qquad
 \mathcal A_h(T[n]):=H_n
\]
とする．再構成は常に蓄積標本全体 $K_n$ を用いる．$\mathcal B_h$ は
集合駆動型だが，$\mathcal A_h$ が最終的に安定する文法は提示順序に依存し得る．
主張するのは，標的言語を生成する文法への最終的な構文的安定化だけである
（系~\ref{cor:ilt}）．

\section{生成語型付き精密化と有限証拠による完全性}
\label{sec:typing}

\subsection{標的の精密化と有限証拠}

再構成作用素と逐次学習器はすでに定義済みである．ここでは完全性の証明のためだけに標的側の精密化を導入する．目標は，対応する正例証拠が出現した後に，仮説文法が有限個の型付き標的導出を模倣できることを示すことである．次の標準形は証明上の道具としてのみ用いる．

\begin{definition}[開始記号分離二分標準形]
CFG $G=(V,\Sigma,P,S_0)$ が \emph{開始記号分離二分標準形}（SSBNF）にあるとは，$S_0$ がいかなる右辺にも現れず\footnote{開始記号を文全体の生成入口として内部の非終端から分離し，型付き精密化と $\lambda$ の特別扱いを単純に保つためである．}，すべての開始規則が $S_0\to A$ または $S_0\to\lambda$ の形であり，すべての非開始規則が $A\to a$ または $A\to BC$ の形であることをいう．
\end{definition}

任意の非空文脈自由言語は，標準的な開始記号分離，無用記号除去，開始記号以外での $\lambda$/単位規則除去，および二分化の手続きにより，同値な縮約済み SSBNF 文法をもつ．標的言語が空集合の場合は自明である．

$L\in\Ccf{h}$ に対する縮約済み SSBNF 文法を $G=(V,\Sigma,P,S_0)$ とする．その \emph{完全生成語型付き精密化} は，開始記号 $\widetilde S$ と，$A\in V\setminus\{S_0\}$, $\mu\in M$ に対する非開始記号 $A_\mu$ をもつ．以下，型添字の併置は $M$ における積を表す．すなわち $\mu\nu=\mu\cdot\nu$ とする．規則は次である．
\[
\begin{array}{ll}
\widetilde S\to A_\mu & (S_0\to A\in P,\ \mu\in M),\\
\widetilde S\to\lambda & (S_0\to\lambda\in P),\\
A_{h(a)}\to a & (A\to a\in P),\\
A_{\mu\nu}\to B_\mu C_\nu & (A\to BC\in P,\ \mu,\nu\in M).
\end{array}
\]
\emph{型付き精密化} $\widetilde G=(\widetilde V,\Sigma,\widetilde P,\widetilde S)$ は，この文法の縮約部分とする．したがって，規約として $\widetilde G$ のすべての記号は到達可能かつ生産的である．生成語型の添字は完全性証明の規則~(R3) に必要である．型付けがなければ，標的規則 $A\to BC$ に対して，$A$ の標準生成語 $\omega(A)$ と $B,C$ の標準生成語の連接 $\omega(B)\omega(C)$ が同じ $h$-値をもつとは限らず，必要な規則~(R3) の遷移を保証できない．生成語型で分割することにより，この等式が不変量となる．一方で周囲の左文脈または右文脈には型を付けない．再構成作用素がそのような型を一度も検査しないからである．

\begin{proposition}[型付き持ち上げと生成語不変量]\label{prop:typed-core}
型付き精密化は $L(\widetilde G)=L(G)$ を満たす．さらに，$X=A_\mu$ が $\widetilde G$ の非開始記号であるとき，$X\Rightarrow_{\widetilde G}^*w$ ならば $h(w)=\mu$ である．
\end{proposition}
\begin{proof}
型注釈を消去すれば，任意の型付き導出は $G$ の導出へ写る．逆に，$G$ の成功した導出木の各非開始節点に，その終端生成語の $h$-型を注釈する．上で定めた終端規則と二項規則は，それらの型のボトムアップ伝播則そのものである．したがって注釈付き木は完全精密化の導出であり，成功した導出木に現れるすべての記号は trim 後も残る．不変量は型付き導出についての帰納法から従う．
\end{proof}

$\Sigma$ 上に順序を一つ固定し，それが誘導する \emph{shortlex} 順序を用いる．すなわち短い語を先に置き，同じ長さなら辞書式順序で比較する．$\widetilde G$ の各非開始記号 $X$ に対し，$\omega(X)$ をその shortlex 最小終端生成語とし，これを $X$ の\emph{標準生成語}と呼ぶ．さらに $X$ の終端到達文脈の中から，$|u_X|+|v_X|$ を最小化するものとして $\chi(X)=(u_X,v_X)$ を選び，同値な候補がある場合は $u_X$ の shortlex 順序，次に $v_X$ の shortlex 順序で決める．$\chi(X)$ を $X$ の\emph{標準文脈}と呼ぶ．$\widetilde G$ は縮約済みなので，このような終端到達文脈は存在する．実際，到達可能性により $X$ は到達可能な文形式中に現れ，周囲の非終端記号の生産性により，その左側および右側のすべての記号を終端文字列へ展開できる．命題~\ref{prop:typed-core} から $h(\omega(A_\mu))=\mu$ である．したがって，任意の型付き二項規則 $X=A_{\mu\nu}\to Y=B_\mu\, Z=C_\nu$ に対し，
\[
 h(\omega(X))=\mu\nu=h(\omega(Y))h(\omega(Z))
 =h(\omega(Y)\omega(Z))
\]
が成り立つ．これこそが完全性証明で規則~(R3) により用いられる等式である．

有限証拠集合を
\begin{equation}\label{eq:W}
\begin{split}
\WitnessSet(\widetilde G):={}&
 \{u_X\omega(X)v_X:X\in\widetilde V\setminus\{\widetilde S\}\}\\
&\cup\{u_Xav_X:X\to a\in\widetilde P\}\\
&\cup\{u_X\omega(Y)\omega(Z)v_X:X\to YZ\in\widetilde P\}\\
&\cup\bigl(\{\lambda\}\text{ if }\lambda\in L\bigr)
\end{split}
\end{equation}
により定める．ここに表示したすべての語は $L$ に属し，集合は有限である．特に，二項規則 $X\to YZ$ に対応する証拠語は
\[
\widetilde S\Rightarrow^*u_XXv_X\Rightarrow u_XYZv_X
\Rightarrow^*u_X\omega(Y)\omega(Z)v_X
\]
により正当化される．アンカーと終端規則証拠についても同様である．$\WitnessSet(\widetilde G)$ の要素を \emph{標準正例証拠} と呼ぶ．要素 $u_X\omega(X)v_X$ は $X$ の一つの標準出現を証拠づける $X$ の \emph{アンカー} であり，他の語は終端または二項の \emph{規則証拠} という別役割がある．

\subsection{厳密再構成と Gold 収束}

\begin{theorem}[有限証拠からの完全性]\label{thm:complete}
非空言語 $L$ に対する任意の縮約済み SSBNF 文法を $G$ とし，固定準同型 $h$ に関するその型付き精密化を $\widetilde G$ とする．有限集合 $K$ が $\WitnessSet(\widetilde G)\subseteq K$ を満たすならば，$L\subseteq L(\widehat G(K))$ である．
\end{theorem}

ここが標的型付けが登場する箇所である．学習器とバッチ作用素は第~\ref{sec:algo}節ですでに定義されており，以下の証明は，それらの仮説非終端記号が生成語型付き標的非終端記号を模倣できることを示す．それ以外では，この主張は純粋に組合せ論的である．型付き文法と包含 $\WitnessSet(\widetilde G)\subseteq K$ だけを用い，代入可能性も追加仮定 $K\subseteq L$ も必要としない．

\begin{proof}
各非開始記号 $X$ に対し，$\widehat X:=[\omega(X):u_X,v_X]$ と書く．~\eqref{eq:W} のアンカー語により，$\widehat X$ は存在する．型付き導出 $X\Rightarrow^*w$ について帰納法を行い，$\widehat X\Rightarrow_{\widehat G(K)}^*w$ を証明する．

最初の規則が $X\to a$ であるとする．証拠 $u_Xav_X\in K$ により $[a:u_X,v_X]$ が存在する．$\omega(X)\ne a$ なら，$h(\omega(X))=h(a)$ なので規則~(R3) が $\widehat X$ を $[a:u_X,v_X]$ に変更する．$\omega(X)=a$ ならこの段階は不要である．その後規則~(R4) が $a$ を生成する．

最初の規則が $X=A_{\mu\nu}\to B_\mu C_\nu$ であるとする．$Y:=B_\mu$ および $Z:=C_\nu$ とおく．~\eqref{eq:W} の規則証拠により三つの観測非終端記号が得られ，規則~(R1) を適用して次を得る．
\begin{equation}\label{eq:binary-reconstruction}
 [\omega(Y)\omega(Z):u_X,v_X]
 \to[\omega(Y):u_X,\omega(Z)v_X]
      [\omega(Z):u_X\omega(Y),v_X].
\end{equation}
$h(\omega(X))=\mu\nu=h(\omega(Y)\omega(Z))$ なので，必要なら規則~(R3) により $\widehat X$ を~\eqref{eq:binary-reconstruction} の左辺へ変更する．すでに$\omega(X)=\omega(Y)\omega(Z)$であるならば，この段階は不要である．次に規則~(R2) により二つの子を $\widehat Y$ および $\widehat Z$ へ輸送する．これらはそれぞれ $Y,Z$ の標準文脈によって定まる．二つの部分導出は帰納法の仮定により完了する．

最後に，任意の非空 $w\in L$ は $\widetilde S\to X\Rightarrow^*w$ という持ち上げられた開始導出をもつ．開始規則そのものにより $(\lambda,\lambda)$ は $X$ の終端到達文脈である．その全長は $0$ であり最小なので，標準選択は $\chi(X)=(\lambda,\lambda)$ となる．そのアンカー $\omega(X)$ は $K$ に属するため，規則~(R5) が上記の導出を開始する．$w=\lambda$ の場合は~\eqref{eq:W} と規則~(R5) から直ちに従う．
\end{proof}

\begin{theorem}[厳密有限標本再構成]
\label{thm:reconstruction-fixed-h}
$L\in\Ccf{h}$ を非空とし，$G$ を $L$ の縮約済み SSBNF 文法，$\widetilde G$ をその型付き精密化とする．$\WitnessSet(\widetilde G)\subseteq K\subseteq L$ ならば，$L(\mathcal B_h(K))=L$ である．
\end{theorem}
\begin{proof}
定理~\ref{thm:soundness} と~\ref{thm:complete} を合わせればよい．
\end{proof}

したがって，任意の非空標的表現 $G$ について，集合 $\WitnessSet(\widetilde G)$ は第~\ref{sec:learning}節の意味で，$\mathcal B_h$ に関する $L(G)$ の特徴標本である．

定理~\ref{thm:reconstruction-fixed-h} は $\{\lambda\}$ を含むすべての非空標的を覆う．残る端点 $L=\emptyset$ については $C_L=\emptyset$ とおけば，$\mathcal B_h(\emptyset)$ は開始生成規則をもたない．これにより定理~\ref{thm:main} の項目~(ii) が確立される．

\begin{corollary}[Gold 同定]\label{cor:ilt}
有限モノイド $M$ への固定準同型 $h:\Sigma^*\to M$ に対し，学習器 $\mathcal A_h$ は $\Ccf{h}$ のすべての非空言語を正例データから同定する．より具体的には，任意の縮約済み SSBNF 標的文法 $G$ とその型付き精密化 $\widetilde G$ に対し，$\WitnessSet(\widetilde G)\subseteq K_{n_0}$ となったとき，段階 $n\ge n_0$ の範囲で生じうる仮説変更は高々一回である．
\end{corollary}

\begin{proof}
本論文のテキスト規約の下では，Gold の主張は非空標的を対象とする．$H_0$ は初期空仮説である．$L$ を非空とし，$T$ を $L$ のテキストとする．$w_1\notin L(H_0)=\emptyset$ なので，学習器は段階 $1$ で再構成する．$\WitnessSet(\widetilde G)\subseteq K_{n_0}$ を満たす $n_0$ を選ぶ．ある段階 $n\ge n_0$ で学習器が再構成するならば，定理~\ref{thm:reconstruction-fixed-h} により新たな仮説はちょうど $L$ となり，その後はいかなる正例も追加の変更を引き起こさない．この場合はこれが最後の仮説変更となる．一方，どの段階 $n\ge n_0$ でも再構成が起こらないと仮定する．健全性から $L(H_{n_0})\subseteq L$ である．さらに，すべての $n$ について $K_n\subseteq L(H_n)$ が成り立つ．再構成時には標本整合性があり，再構成しない段階では新たに提示されたデータがすでに変更されない仮説によって生成されるからである．したがって真の包含 $L(H_{n_0})\subsetneq L$ ならば，段階 $n_0$ までまだ現れていないある $w^\star\in L\setminus L(H_{n_0})$ が存在する．しかし標的のすべての語は最終的にテキストに現れるので，$w^\star$ は後の段階で現れて再構成を強制し，矛盾する．したがって $L(H_{n_0})=L$ である．この場合，段階 $n_0$ 以後に仮説変更は必要ない．
\end{proof}

\section{再構成ステップの計算量}
\label{sec:complexity}

\begin{theorem}\label{thm:poly-build}
$h$ を固定すると，$\mathcal B_h(K)$ は $\|K\|$ に関して多項式時間で構成できる．
\end{theorem}
\begin{proof}
$n_K:=\|K\|$ とおく．標本語には $x\ne\lambda$ を満たす分解 $uxv$ が
$O(n_K^2)$ 個あり，従って $|\widehat V(K)|=O(n_K^2)$ である．因子の $h$-値も
固定モノイド中の $O(n_K^2)$ 回の積でキャッシュできる．規則~(R1) の候補は
$O(n_K^3)$，規則~(R2)--(R3) は高々 $O(n_K^4)$ 個である．因子と文脈を
出現識別子で保持すれば，固定 $h$ の下で各検査は定数時間である．各生成規則の
符号化長は $O(n_K)$ なので，文法全体を明示的に検査・出力する時間は
$O(n_K^5)$ で抑えられる．
\end{proof}

\begin{corollary}\label{cor:poly-update}
保守的学習器 $\mathcal A_h$ の各段階は，それまでに見たデータの符号化サイズに
関して多項式時間で処理できる．
\end{corollary}
\begin{proof}
CFG membership は文法と入力サイズに関して多項式時間であり，現在の仮説サイズは
定理~\ref{thm:poly-build} により多項式である．必要な再構成も同定理による．
\end{proof}

補題~\ref{lem:sample-consistency}，定理~\ref{thm:reconstruction-fixed-h}，
系~\ref{cor:ilt} と合わせて，主定理~\ref{thm:main}(i)--(iii) が得られる
（(ii) のデータサイズ上界を除く）．文法サイズだけによる上界は不可能である．
縮約済み SSBNF 文法
$S_0\to A_n$, $A_0\to a$, $A_i\to A_{i-1}A_{i-1}$ は
$O(n)$ 個の生成規則記号しかもたないが，生成言語は $\{a^{2^n}\}$ である．
また任意の singleton 言語は，任意の $h$ に対して $\sim_h$-代入可能である．

\section{一般の固定型付けに対する型付き厚み}
\label{sec:typed-thickness}

任意の固定 $h$ に対しても，生成語型付けそのもののコストを明示的に測れば，有限証拠構成には直接的な定量上界が得られる．$G=(V,\Sigma,P,S_0)$ を縮約済み SSBNF 文法，$\widetilde G$ をその trim 済み生成語型付き精密化とする．
\[
 N:=|V\setminus\{S_0\}|,\qquad
 N_t:=|\widetilde V\setminus\{\widetilde S\}|
\]
とおき，この表現の \emph{型付き厚み}を
\[
 \tau_h^{\mathrm{typ}}(G):=
 \max\Bigl(\{\,|\omega(X)|:
 X\in\widetilde V\setminus\{\widetilde S\}\,\}\cup\{1\}\Bigr)
\]
と定める．したがって $\tau_h^{\mathrm{typ}}(G)$ は，完全生成語型分割と trim の後に，各非開始記号から得られる最短終端生成語の長さを測る．これは表現に依存するパラメータであり，学習器への追加仮定ではない．

\begin{lemma}[型付き厚みによる証拠語上界]
\label{lem:typed-thickness-bound}
上の $G,\widetilde G$ に対し，$\tau:=\tau_h^{\mathrm{typ}}(G)$ とおく．任意の非開始型付き記号 $X$ について
\[
 |u_X|+|v_X|\le (N_t-1)\tau
\]
が成り立ち，任意の $w\in\WitnessSet(\widetilde G)$ について
\[
 |w|\le (N_t+1)\tau
\]
が成り立つ．さらに，$P_{\mathrm{ter}}$ と $P_{\mathrm{bin}}$ をそれぞれ $G$ の非開始終端規則集合と二項規則集合とすると，
\[
 \|\WitnessSet(\widetilde G)\|
 \le
 \bigl(N|M|+|P_{\mathrm{ter}}|
       +|P_{\mathrm{bin}}||M|^2+1\bigr)
 \bigl((N|M|+1)\tau+1\bigr).
\]
従って固定 $h$ の下では，この上界は $|G|$ と $\tau_h^{\mathrm{typ}}(G)$ の多項式である．
\end{lemma}

\begin{proof}
$\widetilde G$ の非開始記号を頂点とし，各二項生成規則について親から二つの子への辺を張る有向依存グラフを考える．$\widetilde G$ は縮約済みなので，任意の非開始記号 $X$ は到達可能である．従って，ある開始規則の非開始子 $A$ から $X$ への道がこの依存グラフに存在する．そのような道のうち最短のものを選ぶ．最短道は型付き記号を繰り返さないので，辺数は高々 $N_t-1$ である．

道の各辺について，その辺を実現する二項規則を一つ選び，道上にない兄弟記号をその標準生成語へ展開する．$\widetilde S\to A$ から出発してこれらの展開を行うと，$X$ の終端到達文脈が得られる．道外の各標準生成語の長さは高々 $\tau$ なので，文脈の全長は高々 $(N_t-1)\tau$ である．標準文脈 $(u_X,v_X)$ はこれより長くない．

アンカーは標準文脈に標準生成語を一つ，終端規則証拠は終端記号を一つ，二項規則証拠は標準生成語を二つ加える．$\tau\ge1$ だから，式~\eqref{eq:W} の各語の長さは高々 $(N_t+1)\tau$ である．また $N_t\le N|M|$ であり，元の終端規則は高々一つの型付き終端規則を，元の二項規則は高々 $|M|^2$ 個の型付きコピーを生む．従って表示した第一因子は，必要なら含まれる $\lambda$ も含めた証拠語個数を上から抑え，第二因子は各証拠語の $|w|+1$ を上から抑える．両者を掛ければ主張が従う．
\end{proof}

\begin{corollary}[型付き厚みによる特徴データ]
\label{cor:typed-thickness-data}
$h$ を固定する．$G$ が非空な $L(G)\in\Ccf{h}$ に対する縮約済み SSBNF 文法ならば，$\WitnessSet(\widetilde G)$ は $\mathcal B_h$ に対する特徴標本であり，その符号化サイズは $|G|$ と $\tau_h^{\mathrm{typ}}(G)$ の多項式で抑えられる．
\end{corollary}

\begin{proof}
補題~\ref{lem:typed-thickness-bound} が定量上界を与え，定理~\ref{thm:reconstruction-fixed-h} が特徴標本であることを与える．
\end{proof}

しかし，完全型付き精密化を経由する限り，この型付き厚みを元の文法の通常の厚みで一般に置き換えることはできない．

\begin{proposition}[元の厚みと型付き厚みの指数的ギャップ]
\label{prop:typed-thickness-gap}
固定された二元モノイドへの準同型 $h_{\mathtt c}$ と，生成規則中の記号数が $O(n)$ で通常の厚みが $\tau_{G_n^{\mathtt c}}=1$ である縮約済み SSBNF 文法族 $G_n^{\mathtt c}$ が存在し，
\[
 L(G_n^{\mathtt c})=\{\mathtt a,\mathtt c\}^+\in\Ccf{h_{\mathtt c}},
 \qquad
 \tau_{h_{\mathtt c}}^{\mathrm{typ}}(G_n^{\mathtt c})\ge2^n
\]
が成り立つ．従ってこの固定 $h_{\mathtt c}$ に対して，型付き厚みを元の文法サイズと通常の厚みの多項式で一様に抑えることはできない．
\end{proposition}

\begin{proof}
$M_{\mathtt c}=\{1,\zeta\}$ とし，$1$ を単位元，$\zeta$ を零元とする．
\[
 h_{\mathtt c}(\mathtt a)=1,\qquad
 h_{\mathtt c}(\mathtt c)=\zeta
\]
と定める．$n\ge1$ に対し，$G_n^{\mathtt c}$ を次の規則をもつ文法とする．
\[
\begin{aligned}
 S_0&\to Y,\\
 Y&\to \mathtt a\mid\mathtt c\mid YY\mid A_nC,\\
 C&\to\mathtt c,\qquad A_0\to\mathtt a,\\
 A_i&\to A_{i-1}A_{i-1}\mid\mathtt c
 \qquad(1\le i\le n).
\end{aligned}
\]
これは縮約済み SSBNF 文法である．$Y\to\mathtt a\mid\mathtt c\mid YY$ だけですでに $\{\mathtt a,\mathtt c\}$ 上のすべての非空語を生成し，追加規則 $Y\to A_nC$ が生成する語もすべてその集合に属する．従って $L(G_n^{\mathtt c})=\{\mathtt a,\mathtt c\}^+$ である．すべての非終端記号が一文字の終端生成語をもつので $\tau_{G_n^{\mathtt c}}=1$ であり，生成規則中の記号数は $O(n)$ である．

次に，trim 済み完全 $h_{\mathtt c}$-型付き精密化を考える．$\zeta$ は零元なので，$M_{\mathtt c}$ において積が $1$ になるのは $1\cdot1$ だけである．従って
\[
 (A_0)_1\to\mathtt a,\qquad
 (A_i)_1\to(A_{i-1})_1(A_{i-1})_1
 \quad(1\le i\le n)
\]
であり，$(A_i)_1$ から $\mathtt c$ への終端規則は存在しない．帰納法により
\[
 L_{\widetilde G_n^{\mathtt c}}((A_i)_1)=\{\mathtt a^{2^i}\}
\]
を得る．さらに $(A_n)_1$ は trim 後も残る．実際，完全精密化には成功導出
\[
 \widetilde S\to Y_\zeta
 \to (A_n)_1C_\zeta
 \Rightarrow^*\mathtt a^{2^n}\mathtt c
\]
が存在する．従って $(A_n)_1$ の標準生成語の長さは $2^n$ であり，表示した下界が従う．最後に，$\{\mathtt a,\mathtt c\}^+$ では任意の非空因子の分布がすべて $\Sigma^*\times\Sigma^*$ に等しいので，この言語は $\sim_{h_{\mathtt c}}$-代入可能である．また，本論文で固定した符号化の下で $G_n^{\mathtt c}$ の符号化サイズは $O(n)$ の生成規則記号数の多項式であるから，$\tau_{G_n^{\mathtt c}}=1$ であるにもかかわらず $2^n$ は $|G_n^{\mathtt c}|$ に関して超多項式である．
\end{proof}

命題~\ref{prop:typed-thickness-gap} は表現レベルの障害であり，その標的言語自体は短い特徴データをもつ．次の結果は，再構成作用素の内部構成に依存しない通常厚み下界を与える．

\subsection{集合駆動型特徴標本に対する通常厚み下界}
\label{subsec:ordinary-thickness-lower}

二つの単純な観察を用いる．

\begin{lemma}[型ファイバーへの制限]
\label{lem:fibre-restriction}
$h:\Sigma^*\to M$ をモノイド準同型とし，$L$ を $\sim_h$-代入可能とする．
任意の $F\subseteq M$ に対し，$L\cap h^{-1}(F)$ も
$\sim_h$-代入可能である．
\end{lemma}

\begin{proof}
$L_F:=L\cap h^{-1}(F)$ とする．同型の非空因子 $x,y$ が $L_F$ で
共通文脈をもてば，$L$ でも共通文脈をもつので $\D_L(x)=\D_L(y)$ である．
任意の $u,v$ について
\[
 uxv\in L_F
 \iff uxv\in L\ \text{かつ}\ h(u)h(x)h(v)\in F
\]
であり，$x$ を $y$ に置き換えても二条件とも変わらない．
\end{proof}

\begin{lemma}[入れ子標的に対する特徴標本障害]
\label{lem:nested-characteristic-obstruction}
$\mathcal B$ を任意の集合駆動型作用素，$L'\subsetneq L$ とする．
同じ $\mathcal B$ に対し $C,C'$ がそれぞれ $L,L'$ の特徴標本ならば
\[
 C\cap(L\setminus L')\ne\varnothing
\]
である．
\end{lemma}

\begin{proof}
逆に $C\subseteq L'$ とすると，$K:=C\cup C'$ は
$C\subseteq K\subseteq L$ と $C'\subseteq K\subseteq L'$ を同時に満たす．
従って特徴標本性から
$L(\mathcal B(K))=L=L'$ が必要となり，矛盾する．
\end{proof}

アルファベット
\[
 \Sigma_\circ:=\{\mathtt l,\mathtt r,\mathtt a,\mathtt c,
                    \mathtt 0,\mathtt d\}
\]
と，単位元 $1$，零元 $\zeta$ をもつ二元モノイド
$M_\circ:=\{1,\zeta\}$ を固定する．
$h_\circ(\mathtt c)=\zeta$ とし，他の文字をすべて $1$ へ送る準同型
$h_\circ:\Sigma_\circ^*\to M_\circ$ を定める．従って
$h_\circ(w)=\zeta$ は $w$ が $\mathtt c$ を含むことと同値である．

$i\ge0$ に対し
$s_i:=\mathtt c\mathtt0^{\,i+1}\mathtt d$ とおき，
\[
 \mathcal T_0:=\{\mathtt l\mathtt a\mathtt r,,
                    \mathtt l s_0\mathtt r\},
 \qquad
 \mathcal T_i:=\{\mathtt l s_i\mathtt r\}
 \cup\{\mathtt lxy\mathtt r:x,y\in\mathcal T_{i-1}\}
 \quad(i\ge1)
\]
と定める．さらに
\[
 z_0:=\mathtt l\mathtt a\mathtt r,qquad
 z_i:=\mathtt l z_{i-1}z_{i-1}\mathtt r
\]
とおく．$z_i$ は $\mathcal T_i$ の唯一の $\mathtt c$-自由語であり，
\begin{equation}
 |z_i|=5\cdot2^i-2.
 \label{eq:clean-tree-length}
\end{equation}

\begin{lemma}[レベル符号化木言語の代入可能性]
\label{lem:level-coded-substitutable}
任意の $n\ge0$ に対し，$\mathcal T_n$ は Clark--Eyraud の意味で
代入可能である．
\end{lemma}

\begin{proof}
$\mathtt l,\mathtt r$ を木の対応する括弧とみなす．
$\mathcal T_n$ の語は残余高さ $n$ の順序付き二分木を符号化する．
高さ $0$ の本体は $\mathtt a$ または $s_0$，高さ $i>0$ の本体は
$s_i$ または高さ $i-1$ の二部分木の連接である．
\[
 b_0:=\mathtt a,qquad b_i:=z_{i-1}z_{i-1}\quad(i\ge1)
\]
とおくと，任意の語は完全展開木 $z_n$ に対し，互いに祖先・子孫関係にない
節点本体を
\begin{equation}
 s_i\longleftrightarrow b_i
 \tag{$\dagger$}
\end{equation}
で置換して得られる．

この置換位置は語から認識できる．
$s_i=\mathtt c\mathtt0^{i+1}\mathtt d$ の出現は高さ $i$ の短絡本体に限られ，
原始的均衡語 $z_j$ の各出現は高さ $j$ の完全展開部分木である．従って
$z_{i-1}z_{i-1}$ の出現は一つの高さ $i$ 節点の二つの子に一致する．
ゆえに~($\dagger$) の各字面上の出現は合法な本体置換位置である．

いま $uxv,uyv\in\mathcal T_n$ とする．構文解析は決定的である．
二つの解析で短絡／展開の選択が異なる節点本体は，$u$ の直後と $v$ の直前の
二つの切断位置の間に完全に含まれる．実際，本体が左切断より前に始まれば，
最初の本体記号（$\mathtt c$ と $\mathtt a$ または $\mathtt l$）の差が
共通接頭辞内に現れ，右切断より後に終われば，最後の本体記号
（$\mathtt d$ と $\mathtt a$ または $\mathtt r$）の差が共通接尾辞内に現れる．
従って $x$ から $y$ へ，表示因子内部だけで合法な~($\dagger$) の有限置換列が
存在する．任意の $pxq\in\mathcal T_n$ に同じ置換列を適用すれば
$pyq\in\mathcal T_n$ を得る．よって
$\D_{\mathcal T_n}(x)\subseteq\D_{\mathcal T_n}(y)$ であり，逆向きも同様なので
両分布は等しい．
\end{proof}

\begin{theorem}[集合駆動型の通常厚み下界]
\label{thm:ordinary-thickness-lower}
上の固定準同型 $h_\circ$ に対し，各 $n\ge1$ について縮約済み CFG
$R_n,R_n^-$ が存在する．
$\mathcal T_n^-:=\mathcal T_n\setminus\{z_n\}$ とおくと
\[
 L(R_n)=\mathcal T_n,qquad L(R_n^-)=\mathcal T_n^-,
\]
\[
 \mathcal T_n,\mathcal T_n^-\in\Ccf{h_\circ},qquad
 |R_n|+|R_n^-|=O(n),qquad
 \tau_{R_n},\tau_{R_n^-}\le n+5.
\]
$\mathcal B$ を両標的に特徴標本をもつ任意の集合駆動型作用素とする．
このとき $\mathcal T_n$ の任意の特徴標本 $C_n$ は
\[
 \|C_n\|\ge5\cdot2^n-1
\]
を満たす．従って，$\Ccf{h_\circ}$ のすべての標的に特徴標本を与える
集合駆動型作用素には，縮約済み CFG サイズと通常の厚みによる一般の
多項式特徴データ上界は存在しない．特に $\mathcal B_{h_\circ}$ にも適用される．
\end{theorem}

\begin{proof}
$\mathcal T_n$ を生成する縮約済み CFG $R_n$ を
\[
\begin{aligned}
 S&\to A_n,&
 Z_0&\to\mathtt0,&
 Z_i&\to Z_{i-1}\mathtt0 &&(1\le i\le n),\\
 A_0&\to\mathtt l\mathtt a\mathtt r
       \mid\mathtt l\mathtt cZ_0\mathtt d\mathtt r,\\
 A_i&\to\mathtt lA_{i-1}A_{i-1}\mathtt r
       \mid\mathtt l\mathtt cZ_i\mathtt d\mathtt r
       &&(1\le i\le n)
\end{aligned}
\]
で与える．すべての記号は到達可能かつ生産的で，生成規則記号数は $O(n)$ であり，
本論文の固定符号化の下で符号化サイズは $n$ の多項式である．短絡側から $|\omega(A_i)|\le i+5$ が得られるので
$\tau_{R_n}\le n+5$ である．

$R_n^-$ では同じ $Z_0,\ldots,Z_n$ と
$A_0,\ldots,A_{n-1}$ の規則を用い，$S^-\to A_n^-$ を加えて
\[
\begin{aligned}
 A_0^-&\to\mathtt l\mathtt cZ_0\mathtt d\mathtt r,\\
 A_i^-&\to\mathtt l\mathtt cZ_i\mathtt d\mathtt r
 \mid\mathtt lA_{i-1}^-A_{i-1}\mathtt r
 \mid\mathtt lA_{i-1}A_{i-1}^-\mathtt r
 \quad(1\le i\le n)
\end{aligned}
\]
とする．帰納法により $A_i^-$ は $\mathcal T_i$ のうち少なくとも一つの
短絡を含む語をちょうど生成する．従って
$L(R_n^-)=\mathcal T_n^-$ であり，この文法も縮約済みで，生成規則記号数は $O(n)$，符号化サイズは多項式，
厚みは高々 $n+5$ である．

補題~\ref{lem:level-coded-substitutable} により $\mathcal T_n$ は
$\sim_{h_\circ}$-代入可能である．また
\[
 \mathcal T_n^-=\mathcal T_n\cap h_\circ^{-1}(\{\zeta\})
\]
なので，補題~\ref{lem:fibre-restriction} により $\mathcal T_n^-$ も同様である．
両言語は有限なので $\Ccf{h_\circ}$ に属する．

両標的の特徴標本を $C_n,C_n^-$ とする．
$\mathcal T_n\setminus\mathcal T_n^-=\{z_n\}$ なので，
補題~\ref{lem:nested-characteristic-obstruction} から $z_n\in C_n$ が従い，
\[
 \|C_n\|\ge|z_n|+1=5\cdot2^n-1
\]
を得る．一方，二つの文法サイズと通常の厚みは $n$ に関して線形である．
\end{proof}

\paragraph{下界の適用範囲と固定窓上界との整合性．}
下界は大きい方の標的言語単独の複雑さではなく，
$\Ccf{h_\circ}$ という標的クラスに関するものである．
$\mathcal T_n$ 自体は通常の意味で代入可能なので $(0,0)$ 固定窓クラスに属する．
一方 $n\ge1$ では $\mathcal T_n^-$ は通常代入可能ではない．
$z_n=p_n\mathtt a q_n$, $p_n:=\mathtt l^{\,n+1}$ と書き，
$q_n$ 中の別の高さ $0$ 本体 $\mathtt a$ を $s_0$ に置き換えた接尾辞を
$\bar q_n$ とする．すると $p_n\mathtt a$ と $p_ns_0$ は
$\mathcal T_n^-$ で文脈 $(\lambda,\bar q_n)$ を共有するが，
$(\lambda,q_n)$ は $p_ns_0$ の文脈であって $p_n\mathtt a$ の文脈ではない．
両因子の $h_\circ$-型は異なる．

さらに任意の固定 $(k,\ell)$ について，十分大きい $n$ では
$\mathcal T_n^-\notin\Ccf{h_{k,\ell}}$ である．そうでなければ
定理~\ref{thm:reconstruction-fixed-h} により
$\mathcal B_{h_{k,\ell}}$ は二つの入れ子標的の双方に特徴標本をもち，
入れ子標的補題は $\mathcal T_n$ の特徴標本に $z_n$ を強制する．これは，
多項式サイズ，厚み $O(n)$ の $R_n$ に対する
系~\ref{cor:window-transfer} の固定窓多項式上界と，十分大きい $n$ で矛盾する．

一般の下界は第~\ref{sec:learning}節の特徴標本意味論だけを用いる．
$\mathcal B_{h_\circ}$ への適用には定理~\ref{thm:reconstruction-fixed-h} を用いる．
また下界は任意の縮約済み CFG 表現を尺度としており，
Coste--Nicolas~\cite{coste-nicolas2019} の RNF のような制限付き正規形表現に
相対化した上界や，異なる有限標本成功概念を排除しない．従って，
任意の固定 $h$ における定性的 Gold 同定と多項式時間更新から，
通常 CFG 厚みによる多項式特徴データ上界は従わない．

\section{固定窓に対する厚み付きデータ}
\label{sec:fixed-window-thick}

以下の二つの正の定量結果はいずれも，$\tau_h^{\mathrm{typ}}$ を上から抑え，
系~\ref{cor:typed-thickness-data} を適用することで得られる．
$h_{k,\ell}$ では，型付き非終端記号が保存すべき情報は有界長の接頭辞と接尾辞
だけなので，型付き厚みを制御できる．命題~\ref{prop:typed-thickness-gap} は，
一般の $h$ ではこの移送が成り立たないことを示す．表現サイズと厚みに関して
多項式の特徴標本上界は，Eyraud--Heinz--Yoshinaka
~\cite{eyraud-heinz-yoshinaka2016} の用語では identification in polynomial
time and thick data (IPTtD) のデータ部分に対応する．Yoshinaka の定理~1，
補題~6--7 は $(k,\ell)$-代入可能性に対し同型の尺度を与える
~\cite{yoshinaka2008}．ここでは $\mathcal B_{h_{k,\ell}}$ に対して同じ種類の
上界を示すが，同じ正規形や同じ多項式次数は主張しない．

$k,\ell\ge0$ を固定し，$r:=k+\ell$ とする．
$G=(V,\Sigma,P,S_0)$ を縮約済み SSBNF 文法，
$\widetilde G$ をその $h_{k,\ell}$-型付き精密化とし，
\[
 N:=|V\setminus\{S_0\}|,qquad
 N_t:=|\widetilde V\setminus\{\widetilde S\}|,qquad
 \theta_{k,\ell}(G):=
 \begin{cases}
   \tau_G,&r=0,\\[1mm]
   r+(2r-1)N\tau_G,&r>0
 \end{cases}
\]
とおく．

\begin{lemma}[固定窓型付き生成語の有界性]
\label{lem:window-typed-yield}
$\widetilde G$ の任意の非開始型付き記号 $A_\mu$ に対し
\[
 |\omega(A_\mu)|\le\theta_{k,\ell}(G)
\]
である．
\end{lemma}
\begin{proof}
付録~\ref{app:fixed-window-bounds} を参照されたい．
\end{proof}

\begin{lemma}[固定窓文脈と証拠語の有界性]
\label{lem:window-context}
任意の非開始型付き記号 $X$ に対し
\[
 |u_X|+|v_X|\le(N_t-1)\theta_{k,\ell}(G),
\]
また $\WitnessSet(\widetilde G)$ の各語の長さは高々
$(N_t+1)\theta_{k,\ell}(G)$ である．
\end{lemma}
\begin{proof}
補題~\ref{lem:window-typed-yield} と
補題~\ref{lem:typed-thickness-bound} を合わせればよい．
\end{proof}

\begin{theorem}[固定窓に対する多項式厚み付きデータ]
\label{thm:window-thick}
任意に固定した $(k,\ell)$ とアルファベット $\Sigma$ に対し，多項式
$Q_{k,\ell}$ が存在して，$L(G)\in\Ccf{h_{k,\ell}}$ を満たす任意の
縮約済み SSBNF 文法 $G$ について
\[
 \|\WitnessSet(\widetilde G)\|
 \le Q_{k,\ell}(|G|,\tau_G)
\]
が成り立ち，$\WitnessSet(\widetilde G)$ は
$\mathcal B_{h_{k,\ell}}$ に対する特徴標本である．
\end{theorem}
\begin{proof}
補題~\ref{lem:window-typed-yield} から
$\tau_{h_{k,\ell}}^{\mathrm{typ}}(G)\le\theta_{k,\ell}(G)$ であり，
固定 $k,\ell,\Sigma$ の下では右辺は $|G|,\tau_G$ の多項式である．
系~\ref{cor:typed-thickness-data} を適用すればよい．
\end{proof}

上の定理は縮約済み SSBNF 表現について述べる．次の正規化により任意の縮約済み
CFG 表現へ移送できる．

\begin{proposition}[厚みを多項式的に保存する SSBNF 正規化]
\label{prop:thick-ssbnf-normal}
固定多項式 $P_1,P_2$ が存在し，非空言語を生成する任意の縮約済み CFG $R$ を
多項式時間で同値な縮約済み SSBNF 文法 $G$ へ変換でき，
\[
 |G|\le P_1(|R|),\qquad
 \tau_G\le P_2(|R|,\tau_R+1)
\]
が成り立つ．
\end{proposition}

証明は付録~\ref{app:thick-ssbnf} に与える．終端分離と二分化を非開始
$\lambda$-規則除去より先に行うことで，文法サイズと最短非空生成語をともに
多項式に保つ．

\begin{corollary}[固定窓計算量の移送]
\label{cor:window-transfer}
任意の固定 $k,\ell\ge0$ と有限アルファベット $\Sigma$ に対し，
縮約済み CFG $R$ で表される任意の $L\in\Ccf{h_{k,\ell}}$，すなわち
文脈自由な $(k,\ell)$-代入可能標的は，
$\mathcal B_{h_{k,\ell}}$ に対して $|R|$ と $\tau_R$ の多項式サイズの
特徴データをもつ．
\end{corollary}
\begin{proof}
命題~\ref{prop:thick-ssbnf-normal} と定理~\ref{thm:window-thick} を適用する．
同値性は命題~\ref{prop:yl-special} による．
\end{proof}

多項式は固定された窓とアルファベットに依存し，$k,\ell$ を可変とした一様上界は
主張しない．これはバッチ再構成の特徴データに関する主張であり，保守的逐次学習器の
Gold 安定化は系~\ref{cor:ilt} で別途与える．

\section{線形部分クラス}
\label{sec:linear}

線形言語はちょうど 1 ターン・プッシュダウンオートマトンによって受理される
~\cite{ginsburg-spanier1966,autebert1997}．以下の正規化はその文法的対応物を
与え，成功導出に一つの継続する非終端記号の道をもたせる．

\begin{proposition}[linear-spine SSBNF 正規化]\label{prop:linear-normal}
固定多項式 $P_{\mathrm{lin}}$ が存在し，非空言語を生成する任意の線形 CFG $G$ を
多項式時間で同値な縮約済み SSBNF 文法
$G_0=(V_0,\Sigma,P_0,S_0)$ へ変換でき，
$|G_0|\le P_{\mathrm{lin}}(|G|)$ である．さらに各二項生成規則は，
$W_a\to a$ を満たすちょうど一つの新鮮な終端ラッパー子と，
ちょうど一つのそれ以外の非終端子をもつ．
\end{proposition}

構成は付録~\ref{app:linear-normalization} に与える．
$\widetilde G_0$ を型付き精密化，
$n_t:=|\widetilde V_0\setminus\{\widetilde S\}|$ とおく．
各基礎非開始記号は高々 $|M|$ 個の型付きコピーをもつので
$n_t\le|V_0\setminus\{S_0\}|\,|M|$ である．

\begin{lemma}[短い標準生成語]\label{lem:linear-short}
$\widetilde G_0$ の任意の型付き非開始記号 $X$ に対し
\[
 |\omega(X)|\le n_t
\]
である．従って $\tau_h^{\mathrm{typ}}(G_0)\le n_t$ である．
\end{lemma}

\begin{theorem}[線形標的に対する多項式特徴データ]
\label{thm:linear-poly}
$h$ を固定する．線形 CFG で表される $\Clin{h}$ の任意の言語は，
$\mathcal B_h$ に対し，標的表現サイズの多項式で
$\|\cdot\|$-サイズが抑えられる特徴標本をもつ．
\end{theorem}
\begin{proof}
$L=\emptyset$ と $L=\{\lambda\}$ では第~\ref{sec:typing}節の標本を用いる．
それ以外では命題~\ref{prop:linear-normal} の $G_0$ を取る．
補題~\ref{lem:linear-short} から
$\tau_h^{\mathrm{typ}}(G_0)\le n_t\le
|V_0\setminus\{S_0\}|\,|M|$ なので，
系~\ref{cor:typed-thickness-data} により $|G_0|\le P_{\mathrm{lin}}(|G|)$ の
多項式サイズの特徴標本を得る．
\end{proof}

定理~\ref{thm:poly-build} と合わせると，任意の有限標本からのバッチ再構成は
多項式時間であり，選んだ線形 CFG 表現に関して特徴データも多項式である．
ただし $\mathcal B_h$ の出力が線形表現であるとは限らないため，保守的逐次学習器
自体を de~la~Higuera の意味で polynomial time-and-data learner とは呼ばない．

\subsection{非正則な線形 \texorpdfstring{fixed-$h$}{fixed-h} 例}

直前の定理は，正則標的や既知の線形代入可能クラスに限定されない．
次を定める．
\[
\begin{split}
L_{\pm,e}:={}&
\{a^{2m}cb^{2m}:m\ge0\}
\cup
\{a^{2m+1}db^{2m+1}:m\ge0\}\\
&\cup
\{a^neb^n:n\ge0\}.
\end{split}
\]
これは次の線形文法によって生成される．
\[
\begin{array}{lll}
S\to U_{\mathrm{ev}}\mid U_e, &
U_{\mathrm{ev}}\to c\mid aU_{\mathrm{odd}}b, &
U_{\mathrm{odd}}\to d\mid aU_{\mathrm{ev}}b,\\[1mm]
&&U_e\to e\mid aU_eb.
\end{array}
\]

$c,d,e$ を\emph{中心記号}と呼ぶ．
$M_{\pm,e}:=\{1,\kappa_{cd},\kappa_e,0\}$ を，
$1$ が単位元であり，
$\alpha\beta=0$ がすべての
$\alpha,\beta\in\{\kappa_{cd},\kappa_e,0\}$ に対して成り立つ
四元モノイドとする．
すなわち $M_{\pm,e}$ は三元 null semigroup に単位元を付加したものである．
準同型
$h:\{a,b,c,d,e\}^*\to M_{\pm,e}$ を
\begin{equation}\label{eq:lpm-h}
 h(a)=h(b)=1,\qquad
 h(c)=h(d)=\kappa_{cd},\qquad
 h(e)=\kappa_e
\end{equation}
により定める．
従って $h(w)$ は，$w$ が中心記号を含まなければ $1$，
ちょうど一つ含みそれが $c$ または $d$ なら $\kappa_{cd}$，
それが $e$ なら $\kappa_e$，二つ以上含めば $0$ である．

また
$\operatorname{par}(c):=0$,
$\operatorname{par}(d):=1$
とおく．
このとき $m,n\ge0$ と $\zeta\in\{c,d,e\}$ に対して
\begin{equation}\label{eq:lpm-membership}
a^m\zeta b^n\in L_{\pm,e}
\iff
m=n
\ \text{かつ}\
\bigl[
 \zeta=e
 \ \text{または}\
 (\zeta\in\{c,d\}
  \ \text{かつ}\
  m\equiv\operatorname{par}(\zeta)\pmod2)
\bigr].
\end{equation}

\begin{proposition}\label{prop:linear-separator-example}
言語 $L_{\pm,e}$ は線形かつ非正則であり，
\eqref{eq:lpm-h} の準同型に対して $\sim_h$-代入可能である．
一方，Clark--Eyraud 代入可能ではなく，
いかなる固定 $k,\ell\ge0$ に対しても
$(k,\ell)$-代入可能ではない．
\end{proposition}

\begin{proof}
線形性は表示した文法から直ちに従い，
$L_{\pm,e}\cap a^*eb^*=\{a^neb^n:n\ge0\}$ から非正則性も従う．
以下 $L:=L_{\pm,e}$ と書く．

$\sim_h$-代入可能性を示す．同型の非空因子 $x,y$ が $L$ で共通文脈をもつとする．
$L$ の語に現れる因子が型 $0$ をもつことはない．共通型が $1$ なら
$x,y\in a^+\cup b^+$ であり，共通文脈から
$|x|_a-|x|_b=|y|_a-|y|_b$ を得る．この値は $a^+\cup b^+$ の語を一意に
決めるので $x=y$ である．

残る場合は
$x=a^i\zeta b^j$, $y=a^{i'}\zeta'b^{j'}$ と書け，各因子は中心記号を
一つだけ含む．式~\eqref{eq:lpm-membership} より，その文脈は $(a^m,b^n)$
の形に限られ，
\[
\begin{array}{ll}
(a^m,b^n)\in\D_L(a^ieb^j)
 &\Longleftrightarrow n-m=i-j,\\
(a^m,b^n)\in\D_L(a^i\zeta b^j),\ \zeta\in\{c,d\}
 &\Longleftrightarrow n-m=i-j\ \text{かつ}\
 m\equiv\operatorname{par}(\zeta)-i\pmod2.
\end{array}
\]
従って $e$-分岐の分布は $i-j$ により，$c/d$-分岐の分布は
$(i-j,\operatorname{par}(\zeta)-i\bmod2)$ により決まる．
共通文脈は対応する不変量を一致させ，同じ $h$-型であることは両因子が同じ分岐に
属することを保証する．よって $\D_L(x)=\D_L(y)$ である．

一方，$c,e\in L$ は空文脈を共有するが
$(a,b)\in\D_L(e)\setminus\D_L(c)$ なので Clark--Eyraud 代入可能ではない．
さらに $k,\ell$ を固定し，$m:=2\max\{1,k+\ell\}$，
$w_e:=a^keb^\ell$, $w_c:=a^kcb^\ell$ とおく．両因子は同じ長さ $k$ の
接頭辞と長さ $\ell$ の接尾辞をもち，文脈
$(a^{m-k},b^{m-\ell})$ を共有するが，
\[
 (a^{m+1-k},b^{m+1-\ell})
 \in\D_L(w_e)\setminus\D_L(w_c).
\]
従って \hyperref[crit:kl]{固定窓代入可能性の反例基準}より
$L_{\pm,e}\notin\mathsf{KL}$ である．
\end{proof}

従って定理~\ref{thm:linear-poly} は $\mathsf{KL}$ の外にある
非正則線形言語にも適用される．表示した文法は even-linear なので，
Takada の control-set 枠組みからの分離は主張しない．

\section{表現力と境界}
\label{sec:boundaries}

以上の主要な学習結果は，固定された可認識合同関係を対象としていた．ここでは，この枠組みの位置づけに必要な構造的比較だけを記録する．要点は四つある．第一に，可認識合同関係による制御は，正則言語の段階ですでに，任意の固定された接頭辞--接尾辞窓よりも真に表現力が高い．第二に，第~\ref{sec:linear}節で示したように，この真の包含は，非正則な線形言語に対しても成立する．第三に，一般の文脈自由言語に対する定理は，その線形部分クラスだけでは尽くされない．すなわち，ある fixed-$h$ クラスは，固定窓階層全体の外にも位置する非線形な文脈自由言語を含む．第四に，それでもなおクラス $\mathsf{RS}$ は，自然な決定性文脈自由言語のいくつかを含まない．

\subsection{\texorpdfstring{$\mathsf{RS}$}{RS} 内の非線形な標的}

$\Delta:=\{a^nb^n:n\ge0\}$ とおき，その Kleene 閉包 $\Delta^*$ を考える．
語 $z\in\{a,b\}^*$ に対して $\beta(z):=|z|_a-|z|_b$ とし，
$\underline\beta(z)$ を $z$ の全接頭辞（$\lambda$ と $z$ 自身を含む）における
$\beta$ の最小値とする．また，$z_i z_{i+1}=ba$ となる位置 $i$ を
$ba$ の\emph{出現}と呼び，その高さを $\beta(z[1..i])$ とする．

以下の特徴づけを用いる．\textup{($\star$)}
$w\in\Delta^*$ であることと，$\underline\beta(w)\ge0$，$\beta(w)=0$，
かつ $w$ におけるすべての $ba$ の出現が高さ $0$ にあることとは同値である．
実際，$w=a^{n_1}b^{n_1}\cdots a^{n_r}b^{n_r}$ ならば右辺は明らかである．
逆に右辺を仮定し，$w\ne\lambda$ とする．高さが $0$ となる連続する位置で
$w$ を切断すると，各ブロックの内部では高さが正である．従って各ブロックの
内部に $ba$ があれば \textup{($\star$)} の高さ $0$ 条件に反するので，
そのような $ba$ は存在しない．各ブロックはしたがって $a^*b^*$ に属し，
全体のバランスが $0$ なので $a^nb^n$ の形である．$w=\lambda$ の場合は自明である．

語 $w$ の先頭記号と末尾記号を
$\operatorname{fst}(w),\operatorname{lst}(w)\in\{\bot,a,b\}$ とし，
$w=\lambda$ のときはいずれも $\bot$ とする．また
$\operatorname{ba}(w)\in\{0,1\}$ を，$w$ が $ba$ を含むかどうかを表す
フラグとする．$h_\star(w):=(\operatorname{fst}(w),\operatorname{lst}(w),
\operatorname{ba}(w))$ とおく．

連結 $xy$ の先頭記号は $x\ne\lambda$ なら $\operatorname{fst}(x)$，
そうでなければ $\operatorname{fst}(y)$ であり，末尾記号についても対称的である．
また $xy$ が $ba$ を含むのは，$x$ が $ba$ を含むか，$y$ が $ba$ を含むか，
または $\operatorname{lst}(x)=b$ かつ $\operatorname{fst}(y)=a$ の場合に限る．
従って $h_\star(xy)$ は $h_\star(x)$ と $h_\star(y)$ だけから決まる．
ゆえに有限像 $M_\star:=h_\star(\{a,b\}^*)$ 上で
$h_\star(x)\cdot h_\star(y):=h_\star(xy)$ と定めることができる．
この積は良定義であり，全射性と連接の結合律から結合的である．
また $h_\star(\lambda)$ が単位元なので，$M_\star$ は有限モノイドであり，
$h_\star:\{a,b\}^*\to M_\star$ は準同型である．

非空語 $x$ と $d\in\mathbb Z$ に対し，
$d+\underline\beta(x)\ge0$ であり，かつ $x$ 内の各 $ba$ の高さ
$\eta$ について $d+\eta=0$ であるとき，$d$ を $x$ の
\emph{許容入口高さ}と呼ぶ．その集合を $\mathcal E(x)$ と書く．

\begin{proposition}[非線形な fixed-$h_\star$ 標的]
\label{prop:nonlinear-rs-example}
言語 $\Delta^*$ は文脈自由，非正則，かつ非線形であり，
上の有限準同型に対して $\sim_{h_\star}$-代入可能である．さらに，
どの固定 $k,\ell\ge0$ に対しても $(k,\ell)$-代入可能ではない．特に，
$\Delta^*\in\mathsf{RS}_{h_\star}(\{a,b\})\cap\mathrm{CFL}$ かつ
$\Delta^*\in\mathsf{RS}(\{a,b\})\setminus\mathsf{KL}(\{a,b\})$ である．
\end{proposition}

\begin{proof}
文法 $S\to TS\mid\lambda$，$T\to aTb\mid\lambda$ は $\Delta^*$ を生成するので，
$\Delta^*$ は文脈自由である．また $\Delta^*\cap a^*b^*=\Delta$ であり，
$\Delta$ は非正則なので $\Delta^*$ も非正則である．さらに，線形言語は
正則言語との共通部分について閉じている一方，
$\Delta^*\cap a^*b^*a^*b^*=\Delta\Delta$ は線形でない~\cite{autebert1997}．
従って $\Delta^*$ は非線形である．

\emph{$\sim_{h_\star}$-代入可能性．}
非空語 $x,y$ が $h_\star(x)=h_\star(y)$ を満たし，共通文脈 $(u,v)$ をもつとする．
$uxv,uyv\in\Delta^*$ なので全体のバランスを比較すれば $\beta(x)=\beta(y)$ である．

ここで \textup{($\star$)} を $sxt$ に適用すると，固定した文脈 $(s,t)$ に対する
所属判定に必要な情報のうち $x$ に依存するものは，$\operatorname{fst}(x)$，
$\operatorname{lst}(x)$，$\beta(x)$，および $\beta(s)\in\mathcal E(x)$
であるかどうかだけである．実際，$x$ 内部の接頭辞非負性と $ba$ の高さ条件は
$\beta(s)\in\mathcal E(x)$ にちょうど含まれる．境界 $s|x$ の $ba$ 条件は
$\operatorname{fst}(x)$ と $\beta(s)$ で決まり，$x|t$ の条件および $t$ 内部の
すべての高さ条件は，出口高さ $\beta(s)+\beta(x)$ と
$\operatorname{lst}(x)$ によって決まる．従って $\mathcal E(x)=\mathcal E(y)$
を示せば $\D_{\Delta^*}(x)=\D_{\Delta^*}(y)$ が従う．

まず $\operatorname{ba}(x)=\operatorname{ba}(y)=1$ とする．内部に一つでも
$ba$ があれば，その高さ $\eta$ に対して許容入口高さは $d=-\eta$ に固定されるので，
$\mathcal E(x)$ と $\mathcal E(y)$ はいずれも高々一元集合である．共通文脈
$(u,v)$ が存在するため $\beta(u)\in\mathcal E(x)\cap\mathcal E(y)$ であり，
従って $\mathcal E(x)=\{\beta(u)\}=\mathcal E(y)$ である．

次に $\operatorname{ba}(x)=\operatorname{ba}(y)=0$ とする．このとき
$x,y\in a^*b^*$ である．任意の $z\in a^*b^*$ について
$\underline\beta(z)=\min\{0,\beta(z)\}$ なので，
$\mathcal E(z)=\{d\in\mathbb Z:d\ge-\min\{0,\beta(z)\}\}$ である．
$\beta(x)=\beta(y)$ から $\mathcal E(x)=\mathcal E(y)$ が従う．以上により
$\Delta^*$ は $\sim_{h_\star}$-代入可能である．

\emph{固定窓の排除．}
$k,\ell\ge0$ を任意に固定し，$N:=\max\{1,k+\ell\}$ とおく．
$s:=a^Nb^N$，$t:=a^Nb^Na^Nb^N$ とする．両語は長さ
$\max\{1,k+\ell\}$ 以上で，長さ $k$ の接頭辞 $a^k$ と長さ $\ell$ の
接尾辞 $b^\ell$ を共有し，しかも $s,t\in\Delta^*$ なので空文脈を共有する．
一方 $asb=a^{N+1}b^{N+1}\in\Delta^*$ であるが，
$atb=a^{N+1}b^Na^Nb^{N+1}$ では中央の $ba$ が高さ $1$ に現れるため，
\textup{($\star$)} より $atb\notin\Delta^*$ である．従って
$(a,b)\in\D_{\Delta^*}(s)\setminus\D_{\Delta^*}(t)$ であり，
\hyperref[crit:kl]{固定窓代入可能性の反例基準}より $\Delta^*$ は
$(k,\ell)$-代入可能ではない．$k,\ell$ は任意であったから
$\Delta^*\notin\mathsf{KL}(\{a,b\})$ である．
\end{proof}

命題~\ref{prop:linear-separator-example} とこの例は，異なる範囲で
fixed-$h$ 枠組みの広がりを示す．前者は，線形部分クラスに対する多項式特徴データ
定理が，固定窓階層の外にある非正則線形言語にも適用されることを示す．これに対し
$\Delta^*$ は，一般の文脈自由 fixed-$h$ 定理が，固定窓階層の外にある非線形
文脈自由言語にも及ぶことを示す．特に $\Delta^*$ は線形でないため even-linear
でもない．ただし，ここから control-set 法一般とのクラス分離を主張するものではない．

比較のため，固定窓階層の文脈自由部分それ自体も線形言語には含まれないことを記しておく．
$\Gamma:=\{a,b,c,d,e\}$ 上で
\[
 P:=\{a^ncb^n:n\ge0\},\qquad
 L_\times:=PdP
 =\{a^ncb^nd\,a^mcb^m:n,m\ge0\}
\]
とおく．これは Clark--Eyraud が論じた中心記号付き言語族
~\cite{clark-eyraud2007} の端点 $c$ を含む版である．まず $P$ の代入可能性を
直接確認する．$P$ の語の非空因子は，$r>0$ に対する $a^r$，$s>0$ に対する
$b^s$，または $i,j\ge0$ に対する $a^icb^j$ のいずれかである．
二つの因子 $x,y$ が $P$ で共通文脈をもつなら，両方が $c$ を含むか，
どちらも含まないかのいずれかである．一方だけが $c$ を含むなら，もう一方に
唯一の $c$ を供給する共通外側文脈が，前者には二個目の $c$ を作ってしまうからである．
どちらも $c$ を含まない場合，$\beta(z):=|z|_a-|z|_b$ とおけば，
共通文脈で得られる二つの $P$ の語のバランスを比較して
$\beta(x)=\beta(y)$ となる．$c$ を避ける因子は $a$ または $b$ の非空純冪なので，
これは $x=y$ を意味する．両方が $c$ を含む場合，
$x=a^icb^j$，$y=a^{i'}cb^{j'}$ と書く．$x$ の任意の $P$-文脈は
$r+i=s+j$ を満たす $(a^r,b^s)$ の形であり，共通文脈の存在から
$i-j=i'-j'$ を得る．従って同じ方程式が $y$ の $P$-文脈も特徴づけるので，
$\D_P(x)=\D_P(y)$ である．ゆえに $P$ は代入可能である．

さらに $P$ は接頭辞自由かつ接尾辞自由である．実際，$n<m$ なら
$a^ncb^n$ と $a^mcb^m$ の第 $n+1$ 文字はそれぞれ $c$ と $a$ であり，
接尾辞についても対称である．次に $L_\times$ が代入可能であることを示す．
非空因子 $x,y$ が $L_\times$ で共通文脈 $(u,v)$ をもつとする．
$L_\times$ の各語は $d$ をちょうど一つ含むので，$x,y$ が含む $d$ の個数は等しい．
どちらも $d$ を含まない場合，一意な区切り記号は共通外側文脈にある．例えば
$v=rds$ なら $uxr,uyr\in P$ であり，$d\in u$ の場合も対称である．
従って $P$ の代入可能性から $\D_P(x)=\D_P(y)$ を得る．
$x$ の任意の別の $L_\times$-文脈でも，一意な $d$ は二つの外側部分の一方にあり，
対応する $P$-文脈を $x$ から $y$ へ移せる．対称性より
$\D_{L_\times}(x)=\D_{L_\times}(y)$ である．

両因子が $d$ を含む場合は
$x=x_1dx_2$，$y=y_1dy_2$ と書く．このとき
$ux_1,uy_1\in P$ かつ $x_2v,y_2v\in P$ である．
接頭辞自由性と接尾辞自由性から，$i=1,2$ について
$x_i=\lambda$ iff $y_i=\lambda$ が従う．対応する部分が非空なら，
左側では共通文脈 $(u,\lambda)$，右側では $(\lambda,v)$ を用いて
$P$ の代入可能性を適用でき，
$\D_P(x_i)=\D_P(y_i)$ を得る．従って $d$ の左右で文脈を独立に移せるので，
再び $\D_{L_\times}(x)=\D_{L_\times}(y)$ である．
ゆえに $L_\times$ は通常の Clark--Eyraud 型の意味で代入可能であり，
したがって $(0,0)$-固定窓クラス，特に $\mathsf{KL}(\Gamma)$ に属する．また
$S\to XdX$，$X\to aXb\mid c$ により文脈自由である．

最後に，$a,b$ を固定し $c,d,e$ を消去する準同型は $L_\times$ を
\[
 \Delta\Delta=\{a^nb^na^mb^m:n,m\ge0\}
\]
へ写す．もし $L_\times$ が線形なら，線形言語は準同型で閉じているので
$\Delta\Delta$ も線形になるはずである．これは上で用いた
$\Delta\Delta$ の標準的な非線形性~\cite{autebert1997} に反する．従って
$L_\times\in\mathsf{KL}(\Gamma)\cap\mathrm{CFL}$ は非線形である．
同じアルファベット上で命題~\ref{prop:linear-separator-example} は
$\mathsf{KL}(\Gamma)$ の外にある線形言語を与えているので，
$\mathsf{KL}(\Gamma)\cap\mathrm{CFL}$ と $\Gamma$ 上の線形文脈自由言語族は
互いに比較不能である．

\subsection{固定窓階層からの正則分離}

カウンタ・アルファベットを $\Sigma_c:=\{\uparrow,\downarrow,\#\}$ とする．
この小節では $\mathsf{RS}:=\mathsf{RS}(\Sigma_c)$，
$\mathsf{KL}:=\mathsf{KL}(\Sigma_c)$ と略記する．
$w=\delta_1\cdots\delta_n\in\Sigma_c^*$ に対し，$\mathrm{ht}_0(w):=0$ とし，
左から一文字読むごとに $\uparrow$ なら高さを $1$ 増やし，
$\downarrow$ なら $1$ 減らし，$\#$ なら $0$ にリセットすることで
$\mathrm{ht}_i(w)$ を定める．$\rho\ge1$ に対して
$\mathrm{CTR}_\rho:=\{w\in\Sigma_c^*:0\le\mathrm{ht}_i(w)\le\rho
\text{ for all }0\le i\le|w|\}$ とおく．

\begin{proposition}\label{prop:ctr-regular}
任意の $\rho\ge1$ に対して $\mathrm{CTR}_\rho$ は正則であり，
従って $\mathsf{RS}$ に属する．
\end{proposition}
\begin{proof}
状態 $0,1,\ldots,\rho$ で現在の高さを記録し，範囲外に出た後のシンク状態を
一つ加えればよい．$\#$ はすべての非シンク状態を $0$ に戻す．従って
$\mathrm{CTR}_\rho$ は正則であり，命題~\ref{prop:regular-auto} より
$\mathsf{RS}$ に属する．
\end{proof}

\begin{theorem}\label{thm:ctr-non-kl}
任意の $\rho\ge2$ に対して，$\mathrm{CTR}_\rho$ はどの $k,\ell\ge0$ に対しても
$(k,\ell)$-代入可能ではない．
\end{theorem}
\begin{proof}
$k,\ell\ge0$ を固定し，
$s:=(\uparrow\downarrow)^k\uparrow^{\rho-1}\#^\ell$，
$t:=(\uparrow\downarrow)^k\uparrow^\rho\#^\ell$ とおく．$\rho\ge2$ より
$|s|,|t|\ge\max\{1,k+\ell\}$ であり，両語は同じ長さ $k$ の接頭辞と
同じ長さ $\ell$ の接尾辞 $\#^\ell$ をもつ．

$s,t$ はともに $\mathrm{CTR}_\rho$ に属するので，空文脈を共有する．さらに
$\uparrow s$ では，最初の $\uparrow$ の後，$(\uparrow\downarrow)^k$ を高さ $1$
から走るため高さは $1$ と $2$ の間を往復し，その後の $\uparrow^{\rho-1}$
によってちょうど高さ $\rho$ に達する．従って $\rho\ge2$ より
$\uparrow s\in\mathrm{CTR}_\rho$ である．一方 $\uparrow t$ は最初のリセットより
前に高さ $\rho+1$ に達するので $\uparrow t\notin\mathrm{CTR}_\rho$ である．
従って $(\uparrow,\lambda)\in\D_{\mathrm{CTR}_\rho}(s)
\setminus\D_{\mathrm{CTR}_\rho}(t)$ であり，
\hyperref[crit:kl]{固定窓代入可能性の反例基準}より結論が従う．
\end{proof}

\begin{corollary}[正則言語による分離]
$\mathsf{KL}\subsetneq\mathsf{RS}$ であり，この真の包含は正則言語の範囲ですでに生じる．
\end{corollary}
\begin{proof}
各固定 $(k,\ell)$ クラスは命題~\ref{prop:yl-special} により $\mathsf{RS}$ に
含まれる．一方，任意の $\rho\ge2$ について命題~\ref{prop:ctr-regular} と
定理~\ref{thm:ctr-non-kl} から
$\mathrm{CTR}_\rho\in\mathsf{RS}\setminus\mathsf{KL}$ である．
\end{proof}

\subsection{\texorpdfstring{$\mathsf{RS}$}{RS} の外にある自然な決定性文脈自由言語}

次の単純な障害は，以下のカウンタ言語と Dyck 言語が，あらゆる有限モノイド準同型が誘導するクラスの外に落ちる共通理由を切り出すものである．

\begin{lemma}[有限モノイド障害]\label{lem:finite-monoid-obstruction}
$L\subseteq\Sigma^*$ とする．無限集合
$\Xi\subseteq\Sigma^+$ が存在し，互いに異なるすべての $x,y\in\Xi$ に対して
$\D_L(x)\cap\D_L(y)\ne\emptyset$ かつ $\D_L(x)\ne\D_L(y)$ が成り立つとする．
このとき $L\notin\mathsf{RS}(\Sigma)$ である．
\end{lemma}
\begin{proof}
$h:\Sigma^*\to M$ を有限モノイドへの任意の準同型とする．$\Xi$ は無限で $M$ は有限だから，互いに異なる $x,y\in\Xi$ で $h(x)=h(y)$ を満たすものが存在する．その分布は交わるが等しくないので，$\sim_h$-代入可能性に反する．従って任意の有限 $h$ に対して
$L\notin\mathsf{RS}_h(\Sigma)$ である．
\end{proof}

上限を取り除くと，
\[
 \mathrm{CTR}:=
 \{w\in\Sigma_c^*:\mathrm{ht}_i(w)\ge0\text{ がすべての }0\le i\le |w|\text{ について成り立つ}\}
\]
を得る．これは決定性文脈自由言語である．決定性プッシュダウンオートマトンは，$\uparrow$ でスタック記号を一つ積み，$\downarrow$ で一つ取り除き（アンダーフローなら棄却し），$\#$ では入力走査を再開する前に底記号まで取り除く決定的なクリア状態へ移ればよい．

\begin{corollary}[上限なしカウンタの障害]
$\mathrm{CTR}\notin\mathsf{RS}(\Sigma_c)$ である．
\end{corollary}
\begin{proof}
補題~\ref{lem:finite-monoid-obstruction} を
$\Xi=\{\uparrow^i:i\ge1\}$ に適用する．$i<j$ なら，
$(\lambda,\downarrow^i)$ は $\uparrow^i$ と $\uparrow^j$ の双方の文脈である一方，
$(\lambda,\downarrow^j)$ は $\uparrow^j$ の文脈だが $\uparrow^i$ の文脈ではない．
\end{proof}

同じ障害は一種類の括弧の Dyck 言語
\[
 D_1:=\{w\in\{a,b\}^*: |w'|_a\ge|w'|_b
 \text{ が }w\text{ のすべての接頭辞 }w'\text{ について成り立つ},\ 
 |w|_a=|w|_b\}
\]
にも現れる．ここで $|t|_a$ と $|t|_b$ は語 $t$ に現れる $a,b$ の個数を表す．これは標準的な決定性文脈自由言語である~\cite{ginsburg-greibach1966}．

\begin{corollary}\label{cor:dyck-not-rs}
$D_1\notin\mathsf{RS}(\{a,b\})$ である．
\end{corollary}
\begin{proof}
補題~\ref{lem:finite-monoid-obstruction} を
$\Xi=\{b^i a^i:i\ge1\}$ に適用する．$i<j$ のとき，文脈 $(a^j,b^j)$ は
$b^ia^i$ と $b^ja^j$ に共通する一方，$(a^i,b^i)$ は $b^ia^i$ の文脈だが
$b^ja^j$ の文脈ではない．実際，接頭辞 $a^ib^j$ がアンダーフローを起こす．
\end{proof}

この障害は，プッシュダウンの観点から有用な解釈をもつ．興味深い点は，無限に大きくなりうるスタック挙動それ自体が $\mathsf{RS}$ と両立しないわけではない．$L_{\pm,e}$ はすでに
$\mathsf{RS}(\{a,b,c,d,e\})$ に属する非正則 1 ターン言語である．
$\mathrm{CTR}$ と $D_1$ で失敗するのは，より具体的な性質である．すなわち，無限個の因子族が，共通文脈によって互いに比較可能なまま，二つずつ異なる文脈的挙動を担うため，有限な $h$ ではすべての互換でない因子を分離できない．ここではこれを障害としてのみ用い，$\mathsf{RS}\cap\mathrm{CFL}$ の完全なプッシュダウン特徴づけとしては主張しない．

最後の例のために，初等的な閉包性を一つ用いる．
\begin{lemma}[固定語による商]\label{lem:rs-fixed-quotient}
$L$ が $\sim_h$-代入可能で $z\in\Sigma^*$ なら，右商
$L/z:=\{w:wz\in L\}$ も $\sim_h$-代入可能である．
\end{lemma}
\begin{proof}
同じ $h$-型をもつ空でない $x,y$ が $L/z$ において共通文脈 $(u,v)$ をもつなら，
$(u,vz)$ は $L$ における共通文脈であるから $\D_L(x)=\D_L(y)$ である．任意の $(s,t)$ に対して，
\[
 (s,t)\in\D_{L/z}(x)\iff(s,tz)\in\D_L(x)
 \iff(s,tz)\in\D_L(y)\iff(s,t)\in\D_{L/z}(y).
\]
\end{proof}

$S\to aSS\mid b$ が生成する言語は通常 \emph{\L{}ukasiewicz 言語}と呼ばれる．本論文での $D_1$ の符号化では
$L_{\mathrm{Luk}}=D_1b$ を満たす~\cite{autebert1997}．これは標準的な決定性文脈自由言語でもある~\cite{ginsburg-greibach1966}．もし
$L_{\mathrm{Luk}}\in\mathsf{RS}(\{a,b\})$ なら，補題~\ref{lem:rs-fixed-quotient} により
$D_1=L_{\mathrm{Luk}}/b\in\mathsf{RS}(\{a,b\})$ となり，系~\ref{cor:dyck-not-rs} に反する．


\subsection{Clark の congruential family との比較}

Clark~\cite{clark2010congruential} の規約では，
CFG の開始側に単一の開始記号ではなく有限初期集合を許す．
有限初期集合 $I\subseteq V$ に対して
$L_G[I]:=\bigcup_{A\in I}L_G(A)$ と書く．
初期集合 $I$ をもつ CFG $G$ が \emph{congruential} であるとは，
各非終端記号 $A$ について $L_G(A)$ が
$L_G[I]$ の単一の構文合同類に含まれることをいう．
$L_G(A)=\emptyset$ の場合はこの条件を自明に満たす．
有限アルファベット $\Gamma$ に対し，
congruential CFG によって生成される $\Gamma$ 上の言語全体を
$\mathsf{CONG}(\Gamma)$ と書く．

\begin{proposition}[Clark の congruential family との比較]
\label{prop:clark-congruential-comparison}
任意の有限アルファベット $\Gamma$ と有限モノイド準同型
$h:\Gamma^*\to M$ に対して
$\Ccf{h}\subseteq\mathsf{CONG}(\Gamma)$ である．
従って
$\mathsf{RS}(\Gamma)\cap\mathrm{CFL}
 \subseteq\mathsf{CONG}(\Gamma)$
であり，$\Gamma=\{a,b\}$ ではこの包含は真である．
\end{proposition}

\begin{proof}
$L\in\Ccf{h}$ とする．
$L=\emptyset$ なら，生成規則をもたない一つの非終端記号を
初期集合とすればよい．
以下 $L\ne\emptyset$ とする．

$L$ の縮約済み SSBNF 文法を取り，
第~\ref{sec:typing}節の型付き精密化
$\widetilde G=(\widetilde V,\Gamma,\widetilde P,\widetilde S)$
を施す．
非開始型付き非終端記号 $X=A_\mu$ を任意に固定する．
$\widetilde G$ は縮約済みなので，
$X$ にはある終端到達文脈 $(u,v)$ が存在する．
従ってすべての $w\in L_{\widetilde G}(X)$ について
$uwv\in L$ である．
一方命題~\ref{prop:typed-core} より $h(w)=\mu$ であり，
SSBNF には非開始 $\lambda$-生成規則がないため $w\ne\lambda$ である．

そこで $x,y\in L_{\widetilde G}(X)$ を任意にとると，
$x,y$ は同じ $h$-値をもち，かつ共通文脈 $(u,v)$ をもつ非空語である．
$L$ の $\sim_h$-代入可能性から
$\D_L(x)=\D_L(y)$，すなわち $x\equiv_L y$ が従う．
従って，各非開始型付き非終端記号 $X$ について
$L_{\widetilde G}(X)$ は単一の $\equiv_L$-同値類に含まれる．

次に $\widetilde S$ とその生成規則を除き，
$I:=\{X:\widetilde S\to X\in\widetilde P\}$ を初期集合とする．
$\lambda\in L$ なら，新しい非終端記号 $X_\lambda$ と規則
$X_\lambda\to\lambda$ を加え，$X_\lambda$ も $I$ に含める．
$\widetilde S$ はいかなる右辺にも現れないので，
各非開始記号 $X$ の生成語言語はこの操作で変わらない．
また命題~\ref{prop:typed-core} より
$L(\widetilde G)=L$ であるから，得られた文法 $G'$ は
$L_{G'}[I]=L$ を満たす．
各旧非開始記号の生成語言語は上で示したように
一つの $\equiv_L$-同値類に含まれ，
$L_{G'}(X_\lambda)=\{\lambda\}
 \subseteq[\lambda]_{\equiv_L}$ でもある．
従って $G'$ は congruential であり，
$L\in\mathsf{CONG}(\Gamma)$ である．
これにより
$\Ccf{h}\subseteq\mathsf{CONG}(\Gamma)$，
従って
$\mathsf{RS}(\Gamma)\cap\mathrm{CFL}
 \subseteq\mathsf{CONG}(\Gamma)$
が従う．

包含が真であることを示すため，
$\Gamma=\{a,b\}$ 上の一種類の括弧 Dyck 言語 $D_1$ を考える．
任意の $x\in D_1$ と $u,v\in\{a,b\}^*$ に対して
$uxv\in D_1 \Longleftrightarrow uv\in D_1$ である．
実際，$x$ は全体のバランスが $0$ で，
各接頭辞の相対バランスが非負である．
従って $x$ を $u$ と $v$ の間に挿入しても，
挿入後の高さは変化せず，$x$ の内部でも挿入位置の高さを下回らない．
逆に $x$ を削除しても同じ高さ条件は保存される．
ゆえに $x\equiv_{D_1}\lambda$ である．

従って文法
$S\to aSbS\mid\lambda$
は $L_G(S)=D_1\subseteq[\lambda]_{\equiv_{D_1}}$
を満たす congruential grammar であり，
$D_1\in\mathsf{CONG}(\{a,b\})$ である．
一方系~\ref{cor:dyck-not-rs} より
$D_1\notin\mathsf{RS}(\{a,b\})$ である．
従って
$\mathsf{RS}(\{a,b\})\cap\mathrm{CFL}
 \subsetneq\mathsf{CONG}(\{a,b\})$ である．
\end{proof}



\section{結論}

固定有限モノイド型付け $h$ の下で，文脈自由 $\sim_h$-代入可能言語を有限個の
標準正例証拠から厳密再構成でき，保守的学習器がすべての非空標的を多項式時間更新で
Gold 同定することを示した．標準証拠は文法サイズと型付き厚みに関して多項式である．
固定型付け $h_\circ$ では，$\Ccf{h_\circ}$ 全体に特徴標本を与える集合駆動型作用素に，
縮約済み CFG サイズと通常の厚みによる多項式上界は存在しない．一方，任意の固定
接頭辞--接尾辞窓ではその上界が成立し，線形標的では線形文法サイズに関する
多項式上界が成立する．表現力については，固定窓階層を正則・非正則線形・非線形の
各水準で真に拡張する例を与え，$D_1\notin\mathsf{RS}$ および
$\mathsf{RS}(\Gamma)\cap\mathrm{CFL}$ が Clark の congruential family に
真に含まれることを示した．

未解決問題として，$\mathsf{RS}\cap\mathrm{CFL}$ の特徴づけ，有用な型付けの選択，
未知の $h$ の学習，より豊かな文法形式への拡張，および別の表現・有限標本尺度に
基づく効率性が残る．

\section*{データ利用可能性}
本研究では実証データを使用していない．本改訂稿の定理レベルの数学的内容と同期した
Lean 4 形式化は，\texttt{tcs1-v87-formalization-2.0.0} として
Zenodo（DOI: 10.5281/zenodo.23114558）に公開・保存されている．
対応するソースリポジトリは
\texttt{growupkuriyama-hub/tcs1-lean-formalization} である．
本論文中の証明は自己完結しており，Lean 開発は再現性および検証のための
補助的成果物として提供する．

\section*{利益相反の申告}
著者は，本論文で報告した研究に影響を与えたと考えられる既知の競合する金銭的利害関係または個人的関係がないことを申告する．

\section*{研究費}
本研究は，公的機関，民間企業，または非営利部門のいずれの資金配分機関からも特定の助成を受けていない．

\section*{謝辞}
著者は，指導・励まし・関連文献への助力をいただいた金澤誠教授と吉仲亮教授，
ならびに本論文を大幅に改善する丁寧で建設的な報告をくださった二名の匿名査読者に
感謝する．残る誤りはすべて著者の責任である．

\appendix
\section{固定窓証拠語上界と厚み保存正規化}
\label{app:fixed-window-bounds}

\subsection{補題~\ref{lem:window-typed-yield} の証明}

\begin{proof}[補題~\ref{lem:window-typed-yield} の証明]
$r=0$ のとき，$h_{0,0}$ は $\Sigma^+$ 上で定数なので，
生産的な型付き非終端記号 $A_\mu$ は基礎となる非終端記号 $A$ と
同じ非空生成語をもち，$|\omega(A_\mu)|\le\tau_G$ である．

$r>0$ とし，$A_\mu$ が生産的であるとする．
$\mu$ が $M_{k,\ell}$ の短語要素なら，型 $\mu$ の語はその短語自身なので，
その長さは $r$ 未満である．
残る長語要約要素 $\mu=\langle p,q\rangle$ について，
この型をもつ $A$ の導出木を一つ選び，最初の $k$ 個と最後の $\ell$ 個の
終端葉を印付けする．印付け葉はちょうど $r$ 個であり，そのラベルは必要な順序で
$p$ と $q$ を与える．根から印付け葉への経路の和を $T_0$ とする．

$T_0$ 上でちょうど一つの子だけが $T_0$ に残る非終端ノードの極大鎖を考える．
同じ基礎非終端記号が一つの鎖に二度現れたら，
上側の出現を下側の部分木で置き換える．
これは同じ非終端記号からの正しい導出であり，
印付け葉を含まない兄弟部分木しか削除しないので，
印付けられた境界終端記号とその順序は保存される．
これを繰り返すと，各極大一子鎖には高々 $N$ 個の非終端記号しか残らない．

各極大一子鎖を縮約すると，$r$ 個の印付け葉を結ぶ分岐骨格が得られ，
その頂点数は高々 $2r-1$ である．
したがって，最初の分岐頂点より上にあり得る初期鎖も含め，
極大一子鎖は高々 $2r-1$ 本である．
$T_0$ の外側に省略された兄弟部分木が生じるのは
これらの鎖上のノードだけであり，分岐ノードでは両方の子が $T_0$ に含まれる．
従って省略された兄弟部分木は高々 $(2r-1)N$ 個である．

それぞれをその根の最短終端生成語で置き換えると，
一つにつき高々 $\tau_G$ 個の終端記号を寄与する．
得られる語 $w'$ は依然として $A$ から導出され，
接頭辞 $p$ と接尾辞 $q$ を保つので，同じ $h_{k,\ell}$-型 $\mu$ をもち，
\[
|w'|\le r+(2r-1)N\tau_G.
\]

命題~\ref{prop:typed-core} と同じ生成語型注釈により，
この導出は $A_\mu\Rightarrow_{\widetilde G}^*w'$ に持ち上がる．
$A_\mu$ は到達可能であり，この導出木は成功しているので，
木上のすべての型付き記号は trim 後も残る．
したがって $w'$ は $A_\mu$ の型付き生成語であり，
shortlex 最小の型付き生成語 $\omega(A_\mu)$ も同じ長さ上界を満たす．
\end{proof}

\subsection{厚みを多項式的に保存する SSBNF 正規化}
\label{app:thick-ssbnf}

\begin{proof}[命題~\ref{prop:thick-ssbnf-normal} の証明]
$R=(V,\Sigma,P,S)$ を縮約済みで非空言語を生成する文法とし，$n:=|R|$, $\bar\tau:=\tau_R+1$ とおく．文法サイズと最短終端生成語をともに多項式に保つ順序で標準正規化を行う．

まず新しい開始記号 $S_0$ と規則 $S_0\to S$ を導入し，$\lambda\in L(R)$ なら最終的に $S_0\to\lambda$ を残す．長さ 2 以上の右辺に現れるすべての終端記号を新しいラッパー $W_a\to a$ で分離し，長さ 2 を越える右辺を二分化する．得られる文法を $G_{\mathrm{bin}}$ とする．選んだ妥当な符号化の下で，これらの操作の総サイズは $O(n)$ である．元の各非終端記号は同じ終端生成語言語を保つ．新たな各二分化記号は元のある右辺の連続部分を展開し，そこに含まれる元の記号数は高々 $O(n)$ である．その部分に現れる各生産的非終端記号を最短終端生成語で置き換えれば，ある定数 $c_1$ に対して
\[
 \tau_{G_{\mathrm{bin}}}\le c_1 n\bar\tau
\]
が得られる．

次に，$G_{\mathrm{bin}}$ の空語導出可能記号から得られる最短の \emph{非空} 生成語を抑える．$A\Rightarrow_{G_{\mathrm{bin}}}^*w$ を非空終端語への導出とし，導出木の終端葉を一つ選んで根からの経路をたどる．同じ非終端記号が二度現れれば，上側の部分木を下側の部分木で置き換えてサイクルを削除できる．これを繰り返すと経路上の非終端記号数は高々 $|V_{G_{\mathrm{bin}}}|$ となる．各二項ノードで経路外の兄弟記号をその最短終端生成語へ置き換えると，空語導出可能な兄弟記号は空語を寄与してよく，選んだ終端葉が一文字を寄与するので，ある定数 $c_2$ に対して
\[
 \min\{\,|w|>0:A\Rightarrow_{G_{\mathrm{bin}}}^*w\,\}
 \le 1+|V_{G_{\mathrm{bin}}}|\tau_{G_{\mathrm{bin}}}
 \le c_2 n^2\bar\tau
\]
である．

ここで非開始 $\lambda$-規則をすべて除去する．$G_{\mathrm{bin}}$ は二分文法なので，規則 $A\to BC$ から生じるのは，空語導出可能な子に応じて高々 $A\to BC$，$A\to B$，$A\to C$ であり，指数的な空語導出可能記号の部分集合に関する展開は起こらない．通常の構成は各非終端記号の非空終端生成語言語を保存するので，上の長さ上界も保たれる．続いて単位閉包を計算し，非開始単位規則を除去する．各出発非終端記号と，そこから単位規則だけで到達可能な各非単位規則について，その規則を出発非終端記号にコピーする．この段階の前には非終端記号が $O(n)$ 個，非単位規則も $O(n)$ 個なので，生成される規則は高々 $O(n^2)$ 個である．単位規則消去は各非終端記号の言語を保存するため厚み上界を増加させない．最後に無用記号をトリムする．

得られる文法は縮約済みで開始記号が分離され，各非開始規則は終端規則または二項規則であり，$S_0\to\lambda$ が唯一の可能な $\lambda$-規則である．従って SSBNF である．そのサイズは $n$ の多項式，厚みは高々 $c_2n^2(\tau_R+1)$ であり，すべての変換は多項式時間で実行できる．
\end{proof}

\section{linear-spine SSBNF 正規化の証明}
\label{app:linear-normalization}

命題~\ref{prop:linear-normal} の構成は，標準的な開始記号分離，無用記号・空語導出可能規則・単位規則の除去，および二分化を組み合わせ，第~\ref{sec:linear}節で必要な単一 spine 形を保存する．

\begin{proof}[命題~\ref{prop:linear-normal} の証明]
$L(G)=\{\lambda\}$ なら，単一規則 $S_0\to\lambda$ をもつ縮約済み文法を取ればよい．二項生成規則をもたないので，要求される条件は自明に成り立つ．従って以下では，$L(G)$ が空でない語を含むと仮定する．

$S$ を元の開始記号とする．右辺に決して現れない新しい $S_0$ を導入し，$S_0\to S$ を加え，さらに $\lambda\in L(G)$ なら $S_0\to\lambda$ も加える．以後 $S$ は非開始記号として扱う．無用記号を除去し，$\lambda$ を $S_0$ にのみ保存しつつ非開始 $\lambda$-規則を除去し，その後で非開始単位規則を除去する．開始規則には単位規則消去を適用しないので，最終トリミング後に残る各開始規則は，要求どおり $S_0\to A$ または $S_0\to\lambda$ の形になる．この順序により，$\lambda$-消去によって新たに生じた単位規則も除去される．入力文法は線形なので，各右辺には空語導出可能な非終端記号が高々一つしか現れず，各生成規則から追加される代替右辺は高々一つである．
ここで線形性は二箇所で使われる．第一に，空語導出可能規則の展開を多項式に抑える．第二に，消去後も，残る各非単位右辺には非終端記号が高々一つしか現れないことを保証し，これが以下の「ラッパーでない子を一つだけ持つ」分解を可能にする．単位閉包とラッパー鎖による二分化そのものは，一般 CFG に対しても多項式時間の操作である．単位規則消去では有限な単位到達可能閉包を計算する．すなわち，出発非終端記号と，到達可能な非単位生成規則の各組は，高々一つの複製規則を生じる．従って生成規則数も右辺長の総和も多項式に保たれる．

残るのは，新たな単位規則を作らずに二分化することである．以下の二項生成規則の子として現れる各終端記号 $a$ に対して，新しいラッパー $W_a\to a$ を一つ導入する．同値に，すべての $a\in\Sigma$ に対してそのようなラッパーを導入してもよい．
規則 $A\to uBv$ に対して，$u=a_1\cdots a_m$，$v=b_1\cdots b_n$ と書く．非開始単位規則はすでに消去されているので，$m=n=0$ の場合（これは単位規則 $A\to B$ になる）は起こらない．従って $m+n>0$ である．以下，場合分けを明示する．$m=0$ なら必ず $n>0$ である．$n=1$ なら $A\to BW_{b_1}$ を用いる．
$n\ge2$ なら，新しい $R_1,\ldots,R_{n-1}$ を導入し，
\[
 R_1\to BW_{b_1},\qquad
 R_j\to R_{j-1}W_{b_j}\ (2\le j\le n-1),\qquad
 A\to R_{n-1}W_{b_n}
\]
とする．
次に $m>0$ とする．$n=0$ なら $U:=B$ と定める．ここで $B$ はラッパー導入前の文法の非終端記号なので，この二分化で導入した新しい終端ラッパー記号の一つではない．$n>0$ なら新しい $R_1,\ldots,R_n$ を導入して
\[
 R_1\to BW_{b_1},\qquad
 R_j\to R_{j-1}W_{b_j}\quad(2\le j\le n)
\]
とし，$U:=R_n$ とおく．$m=1$ なら単一の二項規則 $A\to W_{a_1}U$ を用いる．$m\ge2$ なら，新しい $J_1,\ldots,J_{m-1}$ を導入して
\[
 A\to W_{a_1}J_1,\qquad
 J_i\to W_{a_{i+1}}J_{i+1}\quad(1\le i\le m-2),\qquad
 J_{m-1}\to W_{a_m}U
\]
とする．
したがって，$m+n>0$ の下で，$m=0$ または $m\ge1$，$n=0$ または $n\ge1$ のすべての場合が覆われる．また最も内側の二項規則は，新しい単位規則ではなく $B$ 自身を含む．

終端記号のみからなる規則 $A\to a_1\cdots a_m$ について，$m=1$ なら $A\to a_1$ をそのまま残す．$m\ge2$ なら，新しい
$\Theta_2,\ldots,\Theta_m$ を導入し，
\[
 A\to W_{a_1}\Theta_2,\qquad
 \Theta_i\to W_{a_i}\Theta_{i+1}\quad(2\le i\le m-1),\qquad
 \Theta_m\to a_m
\]
とする．（$m=2$ のとき，中間の規則族は空である．）従って新たに得られる非開始規則は，終端規則か，あるいはちょうど一つのラッパー子と一つの別の非終端子をもつ二項規則のいずれかである．構成上，ラッパー $W_a$ は二項規則のラッパー子位置にのみ現れ，継続する非終端子としては決して現れない．一方，終端記号のみの構成は，終端規則をもつ新しい非ラッパー鎖記号で終わる．最終トリミングは，この性質，分離開始形，および縮約済み性を保存する．ラッパーと鎖による分解は，すでに多項式に抑えられた右辺長総和に対して線形であるため，全構成のサイズは多項式であり，多項式時間で実行できる．
\end{proof}

\section{短い標準生成語補題の証明}
\label{app:linear-short-proof}

\begin{proof}[補題~\ref{lem:linear-short} の証明]
命題~\ref{prop:linear-normal} の linear-spine SSBNF 形では，各二項 spine step は
$A\to W_aB$ または $A\to BW_a$ の形で，$W_a\Rightarrow a$ だから，
継続する非終端記号の外側にちょうど一つの終端記号を放出する．

$X$ が終端ラッパーなら $|\omega(X)|=1$ である．そうでなければ，
型付き spine 記号 $X$ からの最短終端導出を取り，spine を
$X=X_0,X_1,\ldots,X_t$ と書く．$X_t$ は終端規則を用い，生成語長は $t+1$ である．
もし $X_i=X_j$（$i<j$）なら，二出現の間を削除して，より少ない spine step をもつ
有効な導出を得られるので最短性に反する．従って $t+1$ 個の型付き spine 記号は
相異なり，すべて非開始記号なので $t+1\le n_t$ である．よって
$|\omega(X)|\le n_t$ が従う．
\end{proof}

\begin{thebibliography}{99}

\bibitem{kanazawa1996}
M.~Kanazawa.
\newblock Identification in the limit of categorial grammars.
\newblock \emph{Journal of Logic, Language and Information}, 5(2):115--155, 1996.
\newblock doi:10.1007/BF00173697.


\bibitem{gold1967}
E.\,M.~Gold.
Language identification in the limit.
\textit{Information and Control}, 10(5):447--474, 1967.

\bibitem{pin2025}
J.-E.~Pin.
\newblock \emph{Mathematical Foundations of Automata Theory}.
\newblock Lecture notes, version of March 24, 2025.
\newblock \url{https://www.irif.fr/~jep/PDF/MPRI/MPRI.pdf}.

\bibitem{higuera1997}
C.~de la Higuera.
Characteristic sets for polynomial grammatical inference.
\textit{Machine Learning}, 27:125--138, 1997.

\bibitem{higuera2010}
C.~de la Higuera.
\textit{Grammatical Inference: Learning Automata and Grammars}.
Cambridge University Press, 2010.


\bibitem{clark-eyraud2007}
A.~Clark and R.~Eyraud.
Polynomial identification in the limit of substitutable
context-free languages.
\textit{Journal of Machine Learning Research},
8:1725--1745, 2007.

\bibitem{yoshinaka2008}
R.~Yoshinaka.
Identification in the limit of $k, l$-substitutable
context-free languages.
In \textit{Grammatical Inference: Algorithms and Applications -- ICGI 2008},
LNCS 5278, pp.~266--279, Springer, 2008.
doi:10.1007/978-3-540-88009-7\_21.


\bibitem{takada1988}
Y.~Takada.
\newblock Grammatical inference for even linear languages based on control sets.
\newblock \textit{Information Processing Letters}, 28(4):193--199, 1988.
doi:10.1016/0020-0190(88)90208-6.

\bibitem{takada1995}
Y.~Takada.
\newblock Learning formal languages based on control sets.
\newblock In K.~P. Jantke and S.~Lange (Eds.), \textit{Algorithmic Learning
for Knowledge-Based Systems}, LNAI 961, pp.~316--339, Springer, 1995.
doi:10.1007/3-540-60217-8\_15.

\bibitem{koshiba-makinen-takada1997}
T.~Koshiba, E.~M\"akinen, and Y.~Takada.
\newblock Learning deterministic even linear languages from positive examples.
\newblock \textit{Theoretical Computer Science}, 185(1):63--79, 1997.
doi:10.1016/S0304-3975(97)00016-9.

\bibitem{ginsburg-spanier1966}
S.~Ginsburg and E.~H.~Spanier.
\newblock Finite-turn pushdown automata.
\newblock \textit{SIAM Journal on Control}, 4(3):429--453, 1966.
doi:10.1137/0304034.

\bibitem{coste-typing2004}
F.~Coste, D.~Fredouille, C.~Kermorvant, and C.~de~la~Higuera.
\newblock Introducing domain and typing bias in automata inference.
\newblock In \textit{ICGI 2004}, LNAI 3264, pp.~115--126, Springer, 2004.
doi:10.1007/978-3-540-30195-0\_11.

\bibitem{clark2013}
A.~Clark.
\newblock Learning trees from strings: A strong learning algorithm for some context-free grammars.
\newblock \textit{Journal of Machine Learning Research}, 14:3537--3559, 2013.




\bibitem{eyraud-heinz-yoshinaka2016}
R.~Eyraud, J.~Heinz, and R.~Yoshinaka.
\newblock Efficiency in the identification in the limit learning paradigm.
\newblock In J.~Heinz and J.~M.~Sempere (Eds.),
\emph{Topics in Grammatical Inference}, Springer, 2016.
doi:10.1007/978-3-662-48395-4\_2.

\bibitem{coste-nicolas2019}
F.~Coste and J.~Nicolas.
\newblock Learning local substitutable context-free languages from positive
examples in polynomial time and data by reduction.
\newblock \emph{Proceedings of Machine Learning Research}, 93:155--168, 2019.

\bibitem{autebert1997}
J.-M.~Autebert, J.~Berstel, and L.~Boasson.
\newblock Context-free languages and pushdown automata.
\newblock In G.~Rozenberg and A.~Salomaa (Eds.),
\emph{Handbook of Formal Languages}, Vol.~1, pp.~111--174,
Springer, 1997.

\bibitem{ginsburg-greibach1966}
S.~Ginsburg and S.~A.~Greibach.
\newblock Deterministic context-free languages.
\newblock \emph{Information and Control}, 9:620--648, 1966.

\bibitem{clark2010congruential}
A.~Clark.
\newblock Distributional learning of some context-free languages with a
minimally adequate teacher.
\newblock In \emph{Proceedings of ICGI 2010}, LNCS 6339, pp.~24--37,
Springer, 2010.
doi:10.1007/978-3-642-15488-1\_4.



\end{thebibliography}
\end{document}